\documentclass[12pt,a4paper]{article}

% --------------------------------------------------
% PACOTES BÁSICOS
% --------------------------------------------------
\usepackage[utf8]{inputenc}
\usepackage[T1]{fontenc}
\usepackage[brazil]{babel}
\usepackage{lmodern}
\usepackage[a4paper,margin=2.5cm]{geometry}
\usepackage{setspace}
\usepackage{titlesec}
\usepackage{tocloft}
\usepackage{array}
\usepackage{longtable}
\usepackage{booktabs}
\usepackage{tabularx}
\usepackage{enumitem}
\usepackage{hyperref}
\usepackage{xcolor}
\usepackage{fancyhdr}
\usepackage{lastpage}
\usepackage{graphicx}
\usepackage{nameref}
\usepackage[backend=biber,style=ieee]{biblatex}
\addbibresource{sections/references.bib}

% --------------------------------------------------
% CONFIGURAÇÕES GERAIS
% --------------------------------------------------
\onehalfspacing
\setlength{\parindent}{1.25cm}
\setlength{\parskip}{0.2cm}
\newcommand{\ProjetoNome}{SpinFlow}
\newcommand{\VersaoDocumento}{1.0}
\newcommand{\ProfessoraA}{Professora Heloise Manica Paris Teixeira}
\newcommand{\ProfessoraB}{Raqueline Ritter de Moura Penteado}
\newcommand{\DisciplinaNome}{Engenharia de Software}
\newcommand{\AutorA}{Hélio Toshio Kamakawa R.A. pg56583}
\newcommand{\AutorB}{\textless Nome 2 \textgreater\ \textless R.A. \textgreater}
\newcommand{\AutorC}{Marcelo Figueiredo Terenciani R.A. pg910234}


\hypersetup{
    colorlinks=true,
    linkcolor=black,
    urlcolor=blue,
    citecolor=black,
    pdftitle={Especificação dos Requisitos},
    pdfauthor={Autores do Projeto}
}

% --------------------------------------------------
% CABEÇALHO E RODAPÉ
% --------------------------------------------------
\pagestyle{fancy}
\fancyhf{}
\fancyhead[L]{Especificação dos Requisitos do \ProjetoNome}
\fancyhead[R]{Página \thepage}
\renewcommand{\headrulewidth}{0.4pt}

% --------------------------------------------------
% FORMATAÇÃO DE TÍTULOS
% --------------------------------------------------
\titleformat{\section}
  {\normalfont\bfseries\large}
  {\thesection}{1em}{}

\titleformat{\subsection}
  {\normalfont\bfseries\normalsize}
  {\thesubsection}{1em}{}

\titleformat{\subsubsection}
  {\normalfont\bfseries\normalsize}
  {\thesubsubsection}{1em}{}

\newcommand{\blocoRF}[5]{
\noindent\textbf{#1 \textless #2\textgreater}

\vspace{0.2cm}
\noindent\textbf{Descrição:} #3

\noindent\textbf{Prioridade:} #4

\noindent\textbf{Ator(es):} #5

\noindent\textbf{Requisitos associados:} \textless Ex.: RF002, RF015 \textgreater

\vspace{0.4cm}
}

\newcommand{\blocoNF}[3]{
\noindent\textbf{#1 \textless #2\textgreater}

\vspace{0.2cm}
\noindent\textbf{Descrição:} #3

\vspace{0.4cm}
}

% --------------------------------------------------
% INÍCIO DO DOCUMENTO
% --------------------------------------------------
\begin{document}

% --------------------------------------------------
% CAPA
% --------------------------------------------------
\begin{titlepage}

\noindent\rule{\textwidth}{0.8pt} % traço no topo

\begin{flushright}
    {\Huge \textbf{Especificação dos Requisitos}}
\end{flushright}

\begin{flushright}
    {\Large \textbf{do}}\\[1.0cm]
    
    {\Huge \textbf{\ProjetoNome}}\\[1.2cm]
    
    {\Large \textbf{Versão \VersaoDocumento}}
\end{flushright}

\vspace{1.8cm}
\begin{flushleft}
    \AutorA\\
    \AutorB\\
    \AutorC\\[0.5cm]

    \textbf{Professoras:} \\
        \ProfessoraA\\
        \ProfessoraB\\[0.5cm]
    \textbf{Disciplina:} \DisciplinaNome
\end{flushleft}

\end{titlepage}

% --------------------------------------------------
% SUMÁRIO
% --------------------------------------------------
\tableofcontents
\newpage




% --------------------------------------------------
% REVISÕES
% --------------------------------------------------
\section*{Revisões}
\begin{longtable}{|p{2.5cm}|p{4cm}|p{5cm}|p{3cm}|}
\hline
\textbf{Versão} & \textbf{Autores} & \textbf{Descrição da Versão} & \textbf{Data} \\
\hline
1.0 & Equipe inicial & Criação inicial do documento & \textless dd/mm/aaaa\textgreater \\
\hline
1.1 & Equipe atualizada & Ajustes de conteúdo e refinamento dos requisitos & \textless dd/mm/aaaa\textgreater \\
\hline
\end{longtable}

\newpage

% --------------------------------------------------
% VALIDAÇÕES
% --------------------------------------------------
\section*{Histórico de Validações}
\addcontentsline{toc}{section}{Histórico de Validações}

\section*{Histórico de Validações}
\addcontentsline{toc}{section}{Histórico de Validações}

\begin{longtable}{|>{\raggedright\arraybackslash}p{0.16\textwidth}|>{\raggedright\arraybackslash}p{0.24\textwidth}|>{\raggedright\arraybackslash}p{0.56\textwidth}|}
\hline
\textbf{Data} & \textbf{Com Quem Validou} & \textbf{O Que Validou} \\
\hline
\endfirsthead

\hline
\textbf{Data} & \textbf{Com Quem Validou} & \textbf{O Que Validou} \\
\hline
\endhead

31/03 & \ProfessoraA & Validação da adequação do template IEEE como base estrutural da especificação, incluindo sua coerência com o contexto e os objetivos do projeto. \\
\hline
31/03 & \ProfessoraA & Validação da complexidade do projeto a partir da análise do Diagrama de Classes, considerando aderência ao domínio e viabilidade de implementação. \\
\hline
31/03 & \ProfessoraA & Validação da modelagem dos casos de uso, com ênfase na simplificação aplicada ao Diagrama de Casos de Uso (DCU) para melhoria da legibilidade e compreensão. \\
\hline
31/03 & \ProfessoraA & Validação da utilização da ferramenta Mermaid na elaboração dos diagramas, considerando clareza, consistência visual e adequação à representação dos modelos. \\
\hline

\end{longtable}

\newpage

% --------------------------------------------------
% 1 INTRODUÇÃO
% --------------------------------------------------
% Capitulo 1 da especificacao: introducao, objetivo, escopo, publico-alvo, definicoes, convencoes e referencias.

\section{Introdução}

A \textbf{Pulse Studio Indoor} é uma academia especializada em ciclismo \textit{indoor} e aulas coletivas, cuja operação principal está concentrada nas aulas de \textit{spinning}. Considerando uma estrutura com 1 unidade, 2 salas de \textit{spinning}, 42 \textit{bikes}, 6 professores e cerca de 320 alunos ativos, o contexto do projeto caracteriza uma operação de porte médio, com volume suficiente para exigir maior controle sobre agenda, ocupação das aulas e participação dos alunos.

Nesse cenário, o processo de participação nas aulas deixa de ser trivial. A combinação entre horários de pico, alta procura por vagas, preferência dos alunos por determinadas posições na sala, necessidade de refletir a disponibilidade real das \textit{bikes} e ocorrência de cancelamentos ou manutenções passa a gerar atritos frequentes na rotina da academia. Como consequência, surgem dificuldades de organização do \textit{check-in}, conflitos de reserva, perda de previsibilidade operacional e aumento da insatisfação do aluno diante de falhas de informação ou de inconsistências no acesso às aulas.

Esses problemas tornam-se mais críticos porque, no domínio adotado pelo \ProjetoNome, o aluno não busca apenas uma vaga abstrata. Sua participação depende de elementos concretos como turma, data, mapa da sala, posição, bike e, em certos contextos, o repertório musical associado à aula. Assim, a necessidade central da solução não é apenas registrar reservas, mas sustentar um fluxo de \textit{check-in} confiável, contextualizado e coerente com a operação real da academia.

Dessa forma, o desenvolvimento do \ProjetoNome\ se justifica pela necessidade de organizar o \textit{check-in} das aulas de \textit{spinning} de maneira clara para o aluno e consistente para a operação. A solução proposta busca reduzir atritos de reserva, melhorar a visibilidade sobre disponibilidade e ocupação, preservar a coerência entre mapa, turma e bike e estruturar dados capazes de sustentar histórico, acompanhamento individual e futuras informações gerenciais e estratégicas.

Esta especificação adota como referência de estruturação a norma ISO/IEC/IEEE 29148~\cite{iso29148}, considera as diretrizes aplicáveis de proteção de dados pessoais no contexto da Lei Geral de Proteção de Dados Pessoais (LGPD)~\cite{lgpd} e se apoia no levantamento de requisitos consolidado para o projeto, descrito em Apêndice A (Levantamento de Requisitos).

\subsection{Objetivo do Documento}

Este documento tem como objetivo especificar os requisitos do \ProjetoNome, descrevendo de forma clara e estruturada as funcionalidades esperadas, as restrições relevantes, as regras de negócio que precisam ser respeitadas e os atributos de qualidade que o produto deve atender. A especificação serve como base para desenvolvimento, validação, comunicação entre as partes interessadas e manutenção da solução ao longo de sua evolução.

Além de registrar o comportamento esperado do sistema, este documento formaliza o recorte funcional adotado para o projeto e estabelece um referencial comum para interpretação de termos, prioridades, dependências e limites da solução. Sua organização segue a lógica de uma especificação de requisitos, com ênfase na rastreabilidade entre contexto, escopo, requisitos e artefatos de apoio.

\subsection{Escopo do Produto}

O \ProjetoNome\ é um aplicativo mobile desenvolvido em Flutter, voltado à organização operacional de aulas de \textit{spinning} e centrado no processo de \textit{check-in} do aluno. Seu propósito principal é permitir que a participação em uma aula seja tratada de forma contextualizada, considerando agenda, turma, data, mapa da sala, posição, \textit{bike} e disponibilidade real da operação.

No recorte atual do produto, o escopo funcional está organizado em três frentes de prioridade. A primeira é o \textit{check-in}, que representa o núcleo do sistema e reúne tudo o que é essencial para a reserva funcionar corretamente. A segunda é o suporte ao \textit{check-in}, que inclui elementos que contextualizam ou enriquecem a decisão do aluno, como o \textit{mix} da turma e informações complementares da aula. A terceira é a estrutura complementar do domínio, como grupos de alunos, tratada como prioridade mais baixa e orientada a expansões futuras.

Estão dentro deste escopo os fluxos do aluno relacionados à consulta da agenda, visualização do mapa, reserva por \textit{check-in}, cancelamento, histórico individual e consulta ao repertório da turma. Também estão dentro do escopo os fluxos da professora necessários para manter salas, \textit{bikes}, posições, turmas, manutenções e demais configurações operacionais mínimas que sustentam a reserva e a condução da aula.

Ficam fora do escopo principal desta especificação módulos administrativos amplos, financeiros, comerciais e integrações externas que extrapolem a operação direta da aula e do processo de reserva. Também não constitui foco deste documento a exploração de outras plataformas como frente principal de uso, ainda que a tecnologia adotada permita essa evolução futura. Tais frentes podem compor a evolução do ecossistema da solução, mas não fazem parte do núcleo funcional formalizado neste documento.

\subsection{Público-Alvo}

Este documento destina-se às partes interessadas que precisam compreender, validar, implementar ou avaliar o comportamento esperado do \ProjetoNome. Seu uso principal está associado ao alinhamento entre problema, solução, escopo e requisitos do sistema.

\begin{itemize}[leftmargin=1.2cm]
\item \textbf{Equipe de desenvolvimento}: utilizará a especificação como base para implementação, testes e manutenção do software.
\item \textbf{Gestão da academia}: utilizará o documento para verificar aderência da solução às necessidades operacionais do negócio.
\item \textbf{Profissionais envolvidos na validação}: utilizarão o documento para revisar consistência, completude e coerência entre problema, escopo e requisitos.
\item \textbf{Demais interessados}: poderão consultar o material para compreender objetivos, limites e comportamento esperado do produto.
\end{itemize}

Embora o sistema possua usuários operacionais específicos, como aluno e professora, esses perfis são tratados posteriormente na visão geral do produto. Nesta subseção, o foco está nos leitores da especificação.

\subsection{Definições, Acrônimos e Abreviações}

Esta subseção reúne os principais termos, acrônimos e abreviações utilizados neste documento. O objetivo é estabelecer uma terminologia comum para a leitura dos requisitos, dos modelos de análise e dos artefatos de apoio, reduzindo ambiguidades na interpretação do domínio do sistema.

Os principais termos do domínio adotados nesta especificação são apresentados na Tabela~\ref{tab:glossary}.

\begin{longtable}{|p{4cm}|p{11cm}|}
\caption{Glossário de termos do sistema \ProjetoNome.}\label{tab:glossary} \\ \hline
\textbf{Termo} & \textbf{Descrição} \\ \hline
\endfirsthead

\hline
\textbf{Termo} & \textbf{Descrição} \\ \hline
\endhead

Aluno & Cliente da academia que utiliza o sistema para consultar a agenda, reservar \textit{bike}, acompanhar o histórico de presença e acessar informações do repertório musical das aulas. \\ \hline
Artista ou banda & Entidade cadastral que representa o artista ou a banda responsável por uma música. \\ \hline
\textit{Bike} & Equipamento físico utilizado na aula de \textit{spinning}, associado a um fabricante, a uma posição na sala e, quando aplicável, a registros de manutenção. \\ \hline
Categoria musical & Entidade cadastral utilizada para classificar músicas por categoria, estilo ou finalidade dentro do repertório. \\ \hline
\textit{Check-in} & Registro que representa a reserva do aluno para uma turma em uma data específica, incluindo a posição escolhida na sala e preservando o histórico da participação. \\ \hline
Dia da semana & Enumeração que representa os dias válidos para a recorrência das turmas, de `SEGUNDA' a `DOMINGO'. \\ \hline
Fabricante & Entidade cadastral que identifica a origem da \textit{bike} e concentra seus dados de referência e contato. \\ \hline
Gênero & Enumeração utilizada para representar o gênero declarado no cadastro do aluno, com os valores `MASCULINO', `FEMININO' e `OUTRO'. \\ \hline
Histórico de mix da turma & Registro das associações de mixes a uma turma ao longo do tempo; `dataFim = null' identifica o \textit{mix} ativo em uso naquela turma. \\ \hline
Manutenção & Registro operacional que indica uma intervenção realizada, agendada ou pendente em uma \textit{bike}, podendo afetar sua disponibilidade para reserva. \\ \hline
Mix & Conjunto organizado de músicas que compõe a experiência musical de uma aula ou de um período, com vigência própria e uso efetivo controlado pelo histórico de associação com a turma. \\ \hline
Música & Item do repertório musical que compõe um \textit{mix}, associado a exatamente um artista ou banda e a uma ou mais categorias musicais. \\ \hline
Professora & Responsável operacional que utiliza o sistema para gerenciar salas, \textit{bikes}, posições e turmas, manter dados atualizados para viabilizar o \textit{check-in}, acompanhar a ocupação das aulas, visualizar o mapa operacional e realizar intervenções administrativas conforme as regras do domínio. \\ \hline
Posição da \textit{bike} & Posição ocupável dentro da sala, identificada por fila e coluna, associada a uma \textit{bike} específica e existente apenas na composição da sala. \\ \hline
Sala & Espaço físico onde a turma acontece, definido por uma grade de posições de \textit{bike} organizada por filas e colunas. \\ \hline
Tipo de manutenção & Classificação do problema ou do serviço de manutenção, como defeito mecânico, ajuste ou troca de peça. \\ \hline
Turma & Horário recorrente de aula de \textit{spinning} associado a uma sala, a dias da semana específicos e, quando aplicável, a um \textit{mix} em uso. \\ \hline
Videoaula & Material complementar com orientações sobre a execução de músicas no \textit{spinning}, criado sempre no contexto de pelo menos uma música e podendo servir a mais de uma. \\ \hline

\end{longtable}

\subsection{Convenções}

Este documento adota convenções tipográficas e de identificação com o objetivo de padronizar a apresentação dos requisitos e facilitar sua leitura, rastreabilidade e manutenção. Os requisitos funcionais são identificados pelo prefixo \textbf{RF}, seguido de um identificador numérico sequencial, no formato \textbf{RF001, RF002, RF003, ...}. De forma análoga, os requisitos não funcionais são identificados pelo prefixo \textbf{RNF}, seguido de um identificador numérico sequencial, no formato \textbf{RNF001, RNF002, RNF003, ...}. Os identificadores são únicos e não devem ser reutilizados, mesmo em caso de remoção de requisitos.

As prioridades dos requisitos são classificadas como \textbf{Essencial}, \textbf{Importante} e \textbf{Desejável}, sendo atribuídas individualmente a cada requisito, sem herança automática entre diferentes níveis de especificação. Termos relevantes do domínio, nomes de entidades e identificadores podem ser destacados em \textbf{negrito} quando isso contribuir para a clareza da leitura. Palavras ou expressões em inglês são apresentadas em \textit{itálico} quando apropriado ao contexto.

\subsection{Referências}

Esta subseção reúne as referências utilizadas como base metodológica, normativa e documental para a elaboração desta especificação e para o entendimento do produto.

\nocite{iso29148}
\nocite{lgpd}
\printbibliography[heading=none]


%\subsection{Objetivo do Documento}
%% Especificar o objetivo deste documento e introduzir o produto de software, deixando claro o que será descrito e para quem a especificação foi elaborada.


Este documento tem como objetivo especificar os requisitos do sistema de agendamento e gerenciamento de aulas de \textit{spinning}, descrevendo de forma clara e estruturada suas funcionalidades, restrições e características operacionais. A especificação contempla os requisitos funcionais e não funcionais, bem como as interações entre usuários e o sistema, servindo como base para o desenvolvimento, validação e manutenção do software.

%\subsection{Escopo do Produto}
%% Descrever brevemente o produto e incluir os benefícios esperados, seu propósito principal e quais necessidades ele busca atender.

% Página 67 da ISO 29148

O sistema proposto, denominado \textbf{SpinFlow}, consiste em uma plataforma digital destinada ao apoio de academias na organização e gerenciamento de aulas de \textit{spinning}.

O principal propósito do \textbf{SpinFlow} é melhorar o processo de agendamento e controle das aulas, oferecendo suporte ao gerenciamento de alunos, salas e recursos associados, como bicicletas e estrutura física. Além disso, o sistema contempla funcionalidades auxiliares relacionadas à personalização das aulas, como categorias musicais e artistas, bem como apoio ao controle de manutenção dos equipamentos.

A solução busca atender necessidades recorrentes identificadas no contexto das academias, tais como dificuldades no controle de vagas, organização das turmas, gerenciamento de recursos físicos e ausência de ferramentas estruturadas em ambientes que ainda utilizam processos manuais ou sistemas limitados.

Como benefícios esperados, destacam-se a melhoria na organização das aulas, maior controle sobre a capacidade das salas, otimização do uso dos recursos disponíveis, apoio à gestão de manutenção e aprimoramento da experiência dos alunos.

O escopo do sistema abrange o gerenciamento das seguintes entidades principais: alunos, salas, categorias de música, artistas/bandas, tipos de manutenção e fabricantes, permitindo uma abordagem integrada entre aspectos operacionais e administrativos.

%\subsection{Público-Alvo}
%% Descrever os possíveis leitores do documento, como professor, equipe de desenvolvimento, cliente, usuários finais, avaliadores e demais interessados.

Os possíveis leitores deste documento incluem diferentes partes interessadas envolvidas no desenvolvimento, validação e utilização do sistema \textbf{SpinFlow}. Cada grupo possui objetivos específicos ao consultar esta especificação, desde a implementação técnica até a avaliação e uso do sistema no contexto das academias.

\begin{itemize}[leftmargin=1.2cm]
    \item \textbf{Professor:} responsável pela orientação, acompanhamento e avaliação do documento e da proposta do sistema.
    \item \textbf{Equipe de desenvolvimento:} utilizará a especificação como base para implementação, testes e manutenção do software.
    \item \textbf{Cliente (gestor de academia):} responsável por validar a aderência do sistema às necessidades operacionais da academia e apoiar sua adoção.
    \item \textbf{Usuários finais (instrutores e alunos):} utilizarão o sistema para gerenciamento e participação nas aulas de \textit{spinning}.
    \item \textbf{Avaliadores:} leitores encarregados de analisar a qualidade técnica, a consistência e a completude da especificação.
    \item \textbf{Demais interessados:} outros envolvidos no projeto que necessitem compreender o escopo, os objetivos e as funcionalidades do sistema.
\end{itemize}

%\subsection{Definições, Acrônimos e Abreviações}
%% Definir os termos necessários para compreender o documento.

Esta seção reúne os principais termos, acrônimos e abreviações utilizados neste documento. O objetivo é estabelecer uma terminologia comum para a leitura dos requisitos, dos modelos de análise e das prototipações, reduzindo ambiguidades na interpretação do domínio do sistema.

\paragraph{}Os principais termos do domínio adotados nesta especificação são apresentados na Tabela~\ref{tab:glossary}.

\begin{longtable}{|p{4cm}|p{11cm}|}
    \caption{Glossário de termos do sistema \ProjetoNome.}\label{tab:glossary} \\ \hline
    \textbf{Termo} & \textbf{Descrição} \\ \hline
    \endfirsthead
    \hline
    \textbf{Termo} & \textbf{Descrição} \\ \hline
    \endhead
    Aluno & Cliente da academia que utiliza o sistema para consultar a agenda, reservar bike, acompanhar o histórico de presença e acessar informações do repertório musical das aulas. \\ \hline
    Gênero & Enumeração utilizada para representar o gênero declarado no cadastro do aluno, com os valores `MASCULINO`, `FEMININO` e `OUTRO`. \\ \hline
    Dia da semana & Enumeração que representa os dias válidos para a recorrência das turmas, de `SEGUNDA` a `DOMINGO`. \\ \hline
    Bike & Equipamento físico utilizado na aula de spinning, associado a um fabricante, a uma posição na sala e, quando aplicável, a registros de manutenção. \\ \hline
    Fabricante & Entidade cadastral que identifica a origem da bike e concentra seus dados de referência e contato. \\ \hline
    Tipo de manutenção & Classificação do problema ou do serviço de manutenção, como defeito mecânico, ajuste ou troca de peça. \\ \hline
    Sala & Espaço físico onde a turma acontece, definido por uma grade de posições de bike organizada por filas e colunas. \\ \hline
    Posição da bike & Posição ocupável dentro da sala, identificada por fila e coluna, associada a uma bike específica e existente apenas na composição da sala. \\ \hline
    Turma & Horário recorrente de aula de spinning associado a uma sala, a dias da semana específicos e, quando aplicável, a um mix em uso. \\ \hline
    Check-in & Registro que representa a reserva do aluno para uma turma em uma data específica, incluindo a posição escolhida na sala e preservando o histórico da participação. \\ \hline
    Mix & Conjunto organizado de músicas que compõe a experiência musical de uma aula ou de um período, com vigência própria e uso efetivo controlado pelo histórico de associação com a turma. \\ \hline
    Histórico de mix da turma & Registro das associações de mixes a uma turma ao longo do tempo; `dataFim = null` identifica o mix ativo em uso naquela turma. \\ \hline
    Música & Item do repertório musical que compõe um mix, associado a exatamente um artista ou banda e a uma ou mais categorias musicais. \\ \hline
    Artista ou banda & Entidade cadastral que representa o artista ou a banda responsável por uma música. \\ \hline
    Categoria musical & Entidade cadastral utilizada para classificar músicas por categoria, estilo ou finalidade dentro do repertório. \\ \hline
    Videoaula & Material complementar com orientações sobre a execução de músicas no spinning, criado sempre no contexto de pelo menos uma música e podendo servir a mais de uma. \\ \hline
    Manutenção & Registro operacional que indica uma intervenção realizada, agendada ou pendente em uma bike, podendo afetar sua disponibilidade para reserva. \\ \hline
\end{longtable}

\paragraph{}Os acrônimos e abreviações adotados no documento são apresentados na Tabela~\ref{tab:acronyms}.

\begin{longtable}{|p{4cm}|p{11cm}|}
    \caption{Acrônimos e abreviações utilizados no documento.}\label{tab:acronyms} \\ \hline
    \textbf{Acrônimos / Abreviações} & \textbf{Descrição} \\ \hline
    \endfirsthead
    \hline
    \textbf{Acrônimos / Abreviações} & \textbf{Descrição} \\ \hline
    \endhead
    RF & Requisito Funcional \\ \hline
    RNF & Requisito Não Funcional \\ \hline
    UI & \textit{User Interface}, termo em inglês para Interface do Usuário. \\ \hline
    API & \textit{Application Programming Interface}, termo em inglês para Interface de Programação de Aplicações. \\ \hline
\end{longtable}


%\subsection{Convenções}
%% Definir os termos necessários para compreender o documento.

%\begin{itemize}[leftmargin=1.2cm]
%    \item Requisitos funcionais serão identificados por \textbf{RF001, RF002, ...}
%    \item Requisitos não-funcionais serão identificados por \textbf{NF001, NF002, ...}
%    \item Prioridades adotadas: \textbf{Essencial}, \textbf{Importante} e \textbf{Desejável}
%\end{itemize}

Este documento adota convenções tipográficas e de identificação com o objetivo de padronizar a apresentação dos requisitos e facilitar sua leitura, rastreabilidade e manutenção. Os requisitos funcionais serão identificados pelo prefixo \textbf{RF}, seguido de um identificador numérico sequencial, no formato \textbf{RF001, RF002, RF003, ...}. De forma análoga, os requisitos não funcionais serão identificados pelo prefixo \textbf{RNF}, seguido de um identificador numérico sequencial, no formato \textbf{RNF001, RNF002, RNF003, ...}. Os identificadores são únicos e não devem ser reutilizados, mesmo em caso de remoção de requisitos.

As prioridades dos requisitos serão classificadas como \textbf{Essencial}, \textbf{Importante} e \textbf{Desejável}, sendo atribuídas individualmente a cada requisito, não havendo herança automática entre diferentes níveis de especificação. Além disso, termos relevantes, nomes de entidades e identificadores poderão ser destacados em \textbf{negrito} com o objetivo de facilitar a leitura e a compreensão do documento.

%\subsection{Referências}
%\nocite{*}
%\printbibliography[heading=none]

% --------------------------------------------------
% 2 VISÃO GERAL
% --------------------------------------------------
\section{Visão Geral}

\subsection{Perspectiva do Produto}
Sob a perspectiva desta especificação, o \ProjetoNome\ é um aplicativo \textit{mobile} desenvolvido em Flutter para apoiar a operação de uma academia especializada em \textit{spinning}, com ênfase no \textit{check-in} das aulas e nos elementos que tornam esse fluxo possível. O produto foi concebido para atender à realidade da Pulse Studio Indoor, academia de porte médio cuja rotina combina duas salas de \textit{spinning}, 42 \textit{bikes}, aproximadamente 320 alunos ativos e seis professoras responsáveis pela condução das aulas e manutenção da operação.

Nesse contexto, o produto se posiciona como núcleo operacional do estúdio, conectando agenda, turma, mapa da sala, posição, \textit{bike}, repertório musical e histórico de participação. Sua função não é substituir integralmente todos os controles administrativos da academia, mas oferecer uma solução específica para o domínio da aula de \textit{spinning}, reduzindo atritos de reserva, melhorando a visibilidade sobre a ocupação das aulas e preservando coerência entre os estados operacionais envolvidos. Embora a tecnologia adotada permita evolução para outras plataformas, o foco de uso considerado nesta especificação está em dispositivos móveis, por serem o meio mais aderente à rotina de consulta e reserva do aluno.

Embora o domínio comporte possibilidades de expansão para frentes mais amplas de gestão, relacionamento e inteligência de negócio, o recorte desta especificação permanece concentrado no fluxo do \textit{check-in} e em suas dependências diretas. Isso inclui tanto a experiência do aluno ao reservar uma posição quanto a estrutura operacional mantida pela professora para que a aula aconteça com disponibilidade, organização e consistência.

\subsection{Funcionalidade do Produto}
Em alto nível, o \ProjetoNome\ organiza-se em três frentes de prioridade. A primeira, e principal, é o \textit{check-in}, entendido como o núcleo do produto e como o processo que engloba consulta de agenda, seleção de turma, visualização do mapa da sala, escolha de posição e \textit{bike}, validação da disponibilidade, confirmação da reserva e cancelamento quando permitido. Todos esses elementos existem para viabilizar um \textit{check-in} claro para o aluno e consistente para a operação.

A segunda frente corresponde ao suporte ao \textit{check-in}. Nela se inserem informações e recursos que qualificam a decisão do aluno e enriquecem a experiência da aula, como a consulta ao \textit{mix} associado à turma, ao repertório musical e ao histórico relacionado à participação. Esses elementos não substituem o núcleo operacional, mas ampliam o valor percebido do produto ao contextualizar a reserva.

A terceira frente corresponde às estruturas complementares do domínio, com prioridade inferior no recorte atual. Nessa camada situam-se agrupamentos de alunos e outros elementos de organização que sustentam evolução futura do ecossistema, mas que não definem o valor principal entregue agora. Do ponto de vista operacional da professora, o produto também deve permitir a manutenção das entidades necessárias para o funcionamento do fluxo principal, como salas, \textit{bikes}, posições, turmas, \textit{mixes} e registros de manutenção.

\subsection{Usuários}
Os usuários operacionais principais contemplados nesta especificação são Aluno e Professora. O Aluno representa o perfil central do produto e interage com os fluxos de consulta de agenda, visualização de mapa, escolha de posição, realização de \textit{check-in}, cancelamento de reserva, consulta ao \textit{mix} da turma e acompanhamento do próprio histórico. Sua experiência deve ser direta, rápida e focada na tomada de decisão para participação na aula.

A Professora atua como perfil de suporte operacional ao fluxo do Aluno. Sua responsabilidade está em manter os dados e estados que tornam o \textit{check-in} possível e confiável, incluindo salas, \textit{bikes}, posições, turmas, \textit{mixes}, grupos de alunos e manutenções. Também cabe a esse perfil acompanhar a ocupação das aulas, visualizar o mapa operacional e realizar intervenções administrativas compatíveis com as regras do domínio.

Outros atores, como avaliadores, equipe de desenvolvimento e partes interessadas institucionais, são relevantes para o ciclo do projeto e para o entendimento desta especificação, mas não são tratados como usuários operacionais primários do produto.

\subsection{Ambiente Operacional}
O ambiente operacional previsto para o \ProjetoNome\ prioriza mobilidade, simplicidade de uso e execução local. A aplicação foi estruturada em Flutter, com foco principal em dispositivos móveis utilizados pelos alunos para consulta e reserva de aulas e em dispositivos de apoio da academia utilizados pela professora para manutenção e acompanhamento da operação.

No recorte atual, a solução privilegia armazenamento local como base principal de funcionamento, reduzindo dependência de conectividade contínua e favorecendo a execução dos fluxos essenciais mesmo em cenários de uso offline. Esse comportamento é compatível com a necessidade de resposta rápida em interações como consulta ao mapa, verificação de disponibilidade e confirmação de \textit{check-in}.

Do ponto de vista de uso, o ambiente deve garantir interface legível, tempo de resposta adequado, preservação da integridade dos dados e consistência visual no contexto mobile. A operação esperada não depende, para suas funções centrais, de infraestrutura complexa, serviços externos obrigatórios ou configuração técnica avançada por parte dos usuários finais. A possibilidade de execução em outras plataformas pode ser explorada futuramente, mas não constitui o foco principal do problema tratado neste documento.

\subsection{Restrições de Projeto e Implementação}
O projeto está submetido a restrições técnicas, de domínio e de contexto. Entre as restrições técnicas, destaca-se a necessidade de manter funcionamento prioritariamente offline nas funcionalidades centrais, com persistência local suficiente para sustentar agenda, mapa, reserva, histórico e dados operacionais essenciais. A solução também deve preservar organização modular que permita evolução do domínio sem acoplamento excessivo entre cadastros, regras e interface.

Do ponto de vista do domínio, o produto deve manter coerência entre disponibilidade da \textit{bike}, posição no mapa, estado da turma, manutenção e histórico de \textit{check-in}. Isso implica respeitar regras de unicidade de ocupação, vínculos consistentes entre entidades, tratamento controlado de inativação e manutenção, e preservação de histórico sempre que a regra de negócio exigir rastreabilidade das ações.

Também constituem restrições relevantes o tratamento adequado de dados pessoais no contexto da LGPD, a necessidade de demonstração local da solução e o respeito ao recorte funcional adotado nesta especificação. Dessa forma, módulos administrativos amplos, financeiros ou integrações externas complexas não são tratados como dependências obrigatórias da operação básica documentada neste artefato.

\subsection{Documentação do Usuário}
Considera-se como documentação de apoio ao usuário e às partes interessadas o conjunto de artefatos produzidos ao longo da definição da solução. Isso inclui esta especificação de requisitos, o documento de regras de negócio, os diagramas de casos de uso e de classes, os apêndices com descrições complementares e os protótipos de interface relacionados ao \ProjetoNome.

Também compõem esse conjunto materiais que apoiam entendimento e evolução do produto, como registros estruturados na pasta de apoio do projeto, descrições das prioridades do domínio, documentação complementar de análise e orientações de uso dos fluxos principais do Aluno e da Professora. Em conjunto, esses artefatos oferecem base suficiente para entendimento do sistema, comunicação entre envolvidos e manutenção futura da solução.

\subsection{Suposições e Dependências}
Esta especificação parte da suposição de que a academia disponha de dispositivos aptos a executar a aplicação e de que os usuários principais tenham acesso ao sistema em um contexto operacional controlado. Parte-se também da premissa de que os dados cadastrais fundamentais sejam mantidos com consistência mínima pela Professora, uma vez que o fluxo do Aluno depende diretamente da existência prévia de salas, \textit{bikes}, posições, turmas, \textit{mixes} e demais entidades relacionadas.

Outra dependência central é a disponibilidade do armazenamento local no dispositivo como fonte principal de dados no recorte atual. Além disso, elementos como autenticação robusta, sincronização entre múltiplos dispositivos, integrações externas, notificações e módulos administrativos ampliados devem ser entendidos como possibilidades de evolução do ecossistema, e não como pressupostos obrigatórios para a operação básica aqui especificada.



% \section{Visão Geral}

% \subsection{Perspectiva do Produto}
% Explicar a motivação do desenvolvimento do produto e descrever como ele irá interagir com o ambiente, sistemas externos, usuários e demais componentes envolvidos.

% \subsection{Funcionalidade do Produto}
% Descrever as principais funcionalidades do sistema em alto nível.

% \begin{itemize}[leftmargin=1.2cm]
%     \item Cadastro e autenticação de usuários
%     \item Gerenciamento de dados do domínio do sistema
%     \item Consulta e atualização de informações
%     \item Emissão de relatórios ou visualizações
%     \item Comunicação com serviços externos, se aplicável
% \end{itemize}

% \subsection{Usuários}
% Descrever os perfis de usuários que irão utilizar o sistema.

% \begin{itemize}[leftmargin=1.2cm]
%     \item \textbf{Administrador}: gerencia usuários e configurações
%     \item \textbf{Usuário comum}: utiliza as funcionalidades principais
%     \item \textbf{Operador/Professor/Técnico}: perfil especializado, se houver
% \end{itemize}

% \subsection{Ambiente Operacional}
% Descrever a plataforma, sistema operacional, hardware e demais aplicações que coexistirão com o produto.

% \begin{itemize}[leftmargin=1.2cm]
%     \item Sistema operacional: Windows, Linux, Android, iOS ou Web
%     \item Linguagens e frameworks: conforme projeto
%     \item Banco de dados: conforme projeto
%     \item Navegadores compatíveis: Chrome, Firefox, Edge
% \end{itemize}

% \subsection{Restrições de Projeto e Implementação}
% Descrever limitações técnicas e de projeto.

% \begin{itemize}[leftmargin=1.2cm]
%     \item Restrições de hardware
%     \item Linguagem de programação definida
%     \item Padrões arquiteturais
%     \item Protocolos de comunicação
%     \item Prazo acadêmico de entrega
% \end{itemize}

% \subsection{Documentação do Usuário}
% Listar os documentos de apoio ao usuário.

% \begin{itemize}[leftmargin=1.2cm]
%     \item Manual do usuário
%     \item Guia de instalação
%     \item Guia rápido de uso
%     \item Ajuda embutida no sistema
% \end{itemize}

% \subsection{Suposições e Dependências}
% Descrever as suposições e dependências que podem afetar os requisitos.

% \begin{itemize}[leftmargin=1.2cm]
%     \item Disponibilidade de internet
%     \item Acesso a banco de dados
%     \item Dependência de APIs externas
%     \item Disponibilidade dos usuários para validação
% \end{itemize}

% --------------------------------------------------
% 3 ESPECIFICAÇÃO DAS INTERFACES EXTERNAS
% --------------------------------------------------
\section{Especificação das Interfaces Externas}

\subsection{Requisitos de Interface Externa}

\subsubsection{Interfaces do Usuário}
As interfaces do usuário do \ProjetoNome\ devem atender de forma coerente os dois perfis operacionais contemplados nesta especificação: Aluno e Professora. Como o produto é um aplicativo \textit{mobile} com foco no contexto de uso da academia, a interação deve privilegiar leitura rápida, navegação objetiva e resposta imediata nos fluxos mais frequentes, especialmente consulta de agenda, visualização do mapa, realização de \textit{check-in}, cancelamento de reserva e manutenção operacional.

Para o Aluno, a interface deve priorizar poucos passos, linguagem direta e clareza sobre o estado da aula, da posição escolhida e da disponibilidade real das \textit{bikes}. O mapa da sala deve refletir a organização física do ambiente e distinguir, de maneira compreensível, posições livres, ocupadas, indisponíveis ou afetadas por manutenção. Informações operacionais relevantes ao processo de reserva, como turma, data, horário, \textit{mix} atual e confirmação do \textit{check-in}, devem estar acessíveis de forma simples e contextualizada.

Para a Professora, a interface deve oferecer visão operacional das entidades que sustentam o fluxo principal do produto, incluindo salas, \textit{bikes}, posições, turmas, \textit{mixes}, grupos de alunos e manutenções. Nessa frente, a solução deve favorecer a compreensão da relação entre os elementos do domínio, evidenciar os impactos de ativação, inativação, manutenção e cancelamento e permitir acompanhamento da ocupação das aulas de forma compatível com a rotina da academia.

Independentemente do perfil, a interface não deve depender exclusivamente de cor para representar estados críticos. Situações como disponibilidade, indisponibilidade, bloqueio, ocupação e erro devem possuir reforço textual, icônico ou visual equivalente. O sistema também deve fornecer retorno explícito para ações de sucesso, falha, ausência de dados, conflito de reserva e carregamento de informações.

\subsubsection{Interfaces de Hardware}
No escopo desta especificação, o \ProjetoNome\ não depende de integração direta com hardware especializado, sensores, leitores, catracas, impressoras ou periféricos externos específicos. A interface de hardware relevante é a dos próprios dispositivos utilizados para executar o aplicativo, com destaque para \textit{smartphones} Android e dispositivos iOS utilizados pelos alunos, além de dispositivos móveis ou computadores de apoio utilizados pela Professora na operação da academia.

Os dispositivos de uso principal devem oferecer, no mínimo, tela sensível ao toque com área útil suficiente para leitura da agenda e visualização do mapa da sala, capacidade de armazenamento local para persistência dos dados do aplicativo e recursos de processamento compatíveis com renderização do mapa, consultas e operações de \textit{check-in} sem degradação perceptível do fluxo. O projeto deve ser executável em aparelhos compatíveis com as configurações mínimas de \textit{build} definidas para Android e iOS no ambiente da aplicação, sem exigir acessórios adicionais para que as funções centrais do produto sejam realizadas.

No recorte atual, não há necessidade de câmera, microfone, geolocalização, Bluetooth ou leitura de arquivos locais para execução das funções centrais. Como o produto trabalha com URLs informadas em cadastro para fotos, links e conteúdos externos, a exibição desses recursos depende apenas da capacidade normal de renderização do dispositivo e, quando aplicável, de conectividade para acesso ao endereço remoto correspondente.

\subsubsection{Interfaces de Software}
As interfaces de software do \ProjetoNome\ concentram-se principalmente na comunicação entre a camada de apresentação do aplicativo, as regras de domínio e os mecanismos de persistência local que sustentam o funcionamento do produto. No recorte atual, o sistema deve preservar consistência entre agenda, turma, mapa, posição, \textit{bike}, \textit{mix}, histórico e manutenção, refletindo nas telas exatamente o estado operacional mantido pela aplicação.

Como a solução foi concebida com foco em operação local, a principal interface de software externa ao fluxo visual corresponde ao mecanismo de armazenamento persistente utilizado pelo aplicativo. Essa interface deve permitir gravação, consulta e atualização dos dados do domínio de forma íntegra, especialmente nas operações críticas de \textit{check-in}, cancelamento e atualização da disponibilidade da turma.

No estado atual do projeto, não foram identificadas integrações obrigatórias com serviços externos de autenticação, notificações \textit{push}, captura de mídia ou sincronização remota contínua. Também não foram observadas, nos artefatos móveis principais, permissões sensíveis específicas para câmera, galeria, microfone, localização ou contatos. Assim, a operação básica do produto deve permanecer desvinculada dessas autorizações.

O produto admite, como evolução futura, interfaces complementares com serviços de autenticação, sincronização, notificações e fontes externas de informação. Entretanto, tais integrações não constituem dependências obrigatórias para o funcionamento básico especificado neste documento e não devem comprometer a autonomia dos fluxos centrais do sistema.

\subsubsection{Interfaces de Comunicação}
No recorte atual, o \ProjetoNome\ não depende de comunicação contínua em rede para executar suas funções centrais. Os fluxos principais de consulta, visualização do mapa, reserva, cancelamento e acompanhamento operacional devem permanecer utilizáveis mesmo na ausência de conectividade externa, desde que o contexto local do aplicativo esteja disponível.

As necessidades de comunicação identificadas no projeto estão associadas principalmente ao acesso eventual a endereços externos informados por URL, como imagens remotas de perfil, fotos de artistas, links musicais e referências de vídeo-aula. Esse tipo de comunicação é complementar ao fluxo principal e não deve impedir o uso do aplicativo quando indisponível, podendo apenas limitar a exibição do conteúdo remoto correspondente.

Nos artefatos móveis atualmente configurados, a permissão de internet aparece apenas nos manifestos de desenvolvimento do Android, compatível com execução e depuração durante o ciclo de implementação. No recorte funcional descrito nesta especificação, não há dependência obrigatória de comunicação em rede para reserva, agenda, mapa ou manutenção operacional.

Ainda assim, o contexto da solução admite interfaces de comunicação complementares para frentes futuras do ecossistema, como sincronização entre dispositivos, envio de notificações, integração com serviços remotos e troca de dados com componentes externos da academia. Quando presentes, essas comunicações devem ocorrer de forma controlada, sem alterar a previsibilidade do fluxo principal nem introduzir dependência obrigatória para o uso cotidiano do produto.

Sempre que houver troca de dados entre o aplicativo e componentes externos, a comunicação deverá respeitar as restrições de segurança, privacidade e finalidade aplicáveis ao domínio, preservando a integridade das informações e o tratamento adequado dos dados pessoais no contexto da LGPD.

% \section{Especificação das Interfaces Externas}

% \subsection{Requisitos de Interface Externa}

% \subsubsection{Interfaces do Usuário}
% As interfaces do usuário devem ser intuitivas, responsivas e coerentes com o perfil dos usuários do sistema. Protótipos, wireframes ou telas podem ser entregues em anexo.

% \subsubsection{Interfaces de Hardware}
% Descrever as conexões do produto com hardwares utilizados, quando aplicável, como leitores, sensores, impressoras, dispositivos móveis ou periféricos específicos.

% \subsubsection{Interfaces de Software}
% Descrever as conexões com softwares externos, bibliotecas, frameworks, bancos de dados, serviços em nuvem ou APIs consumidas pelo sistema.

% \subsubsection{Interfaces de Comunicação}
% Descrever como serão realizadas as comunicações, como navegador web, e-mail, formulários eletrônicos, REST API, sockets ou protocolos específicos.

% --------------------------------------------------
% 4 REQUISITOS FUNCIONAIS
% --------------------------------------------------
\section{Requisitos Funcionais}

% Garante numeracao ate o nivel de subsubsection no documento
\setcounter{secnumdepth}{3}

Os requisitos funcionais estão organizados em grupos identificados por siglas que em conjunto formam identificadores únicos. Ator e prioridade são
declarados no nível do grupo e herdados pelos requisitos; exceções são indicadas no próprio requisito. Os critérios de aceite são apresentados
como lista de condições verificáveis.

\subsection{Comportamentos Transversais dos Cadastros}

Para promover consistência e reduzir repetições, os seguintes comportamentos transversais aplicam-se a todos os requisitos de manutenção de cadastros (REQ-CAD-*), salvo indicação explícita em contrário:

\begin{itemize}[leftmargin=0.6cm]
    \item \textbf{Confirmação explícita para operações críticas:} O sistema deve exigir confirmação explícita do usuário para operações de inativação, reativação, cancelamento ou exclusão lógica, quando aplicáveis.
    \item \textbf{Mensagens de erro associadas a campos:} O sistema deve exibir mensagem de erro junto ao campo que causou a rejeição, sem apagar os demais campos do formulário.
    \item \textbf{Preservação de dados em formulários:} O sistema deve preservar os valores informados nos campos do formulário durante validações e tentativas de salvamento, exceto quando explicitamente indicado.
    \item \textbf{Atualização imediata de consultas:} Após operação concluída com sucesso, o sistema deve refletir imediatamente o resultado nas listagens e consultas relevantes.
    \item \textbf{Pré-preenchimento em edição:} O sistema deve apresentar formulários de edição pré-preenchidos com os dados existentes do registro, garantindo consistência e facilitando correções.
    \item \textbf{Impedimento de exclusão física com vínculos:} O sistema deve impedir exclusão física de registros com vínculos associados, exibindo mensagem explicativa.
    \item \textbf{Manutenção de dados após inativação:} O sistema deve manter o cadastro de um registro inativo consultável com seus dados íntegros, sem exclusão física ou ocultação do estado.
    \item \textbf{Seleção controlada de associações:} Para campos que referenciam outras entidades, o sistema deve permitir seleção apenas a partir de opções válidas e ativas apresentadas, impedindo digitação livre e associação a registros inativos em novos vínculos quando aplicável.
    \item \textbf{Imutabilidade do identificador interno:} O sistema deve manter o identificador interno do registro único e inalterável após a criação, preservando-o durante qualquer operação de atualização.
    \item \textbf{Disponibilidade em consultas:} O sistema deve recuperar registros ativos e inativos nas consultas de cadastro, salvo indicação contrária no requisito.
\end{itemize}

Estes comportamentos são herdados implicitamente pelos requisitos de cadastro, sendo omitidos das listas de critérios de aceite individuais para evitar repetição. Exceções ou especificidades são indicadas explicitamente no requisito correspondente.

% Contadores hierárquicos internos para os identificadores RF
% O requisito incorpora a sigla do macrogrupo atual e o numero
% do grupo atual, acrescentando apenas sua sequencia local.
\newcounter{rfmacrogrupo}
\newcounter{rfgrupo}[rfmacrogrupo]
\newcounter{rfreq}[rfgrupo]
\newcounter{rfcadastrotopico}[rfmacrogrupo]
\newcommand{\rfsiglasecao}{RF}
\newcommand{\rfmacrosigla}{}
% \newcommand{\rfgrupoid}{\thesection.\arabic{rfmacrogrupo}.\arabic{rfgrupo}}
\newcommand{\rfgrupoid}{\thesubsection.\arabic{rfgrupo}}
% \newcommand{\rfcadastrotopicoid}{\arabic{rfmacrogrupo}.\arabic{rfcadastrotopico}}
%% \newcommand{\rfcadastrotopicoid}{\thesubsection.\arabic{rfcadastrotopico}}
\newcommand{\rfcadastrotopicoid}{%
  \ifnum\value{rfgrupo}=0
    \thesubsection.\arabic{rfcadastrotopico}%
  \else
    \thesubsection.\arabic{rfgrupo}.\arabic{rfcadastrotopico}%
  \fi
}
\renewcommand{\therfreq}{\rfgrupoid.\arabic{rfreq}}
\newcommand{\rfreqid}{\rfsiglasecao-\rfmacrosigla-\therfreq}

\newcommand{\rfmacroprioridadebase}{}
\newcommand{\rfmacroprioridadepadrao}{}

\makeatletter
\def\reqcad@split#1 #2\relax{%
  \gdef\reqcad@code{#1}%
  \gdef\reqcad@title{#2}%
}
\gdef\reqcad@isreq{0}
\def\reqcad@testREQ#1#2#3#4\@nil{%
  \if R#1%
    \if E#2%
      \if Q#3%
        \gdef\reqcad@isreq{1}%
      \else
        \gdef\reqcad@isreq{0}%
      \fi
    \else
      \if F#2%
        \if -#3%
          \gdef\reqcad@isreq{1}%
        \else
          \gdef\reqcad@isreq{0}%
        \fi
      \else
        \gdef\reqcad@isreq{0}%
      \fi
    \fi
  \else
    \gdef\reqcad@isreq{0}%
  \fi
}
\makeatother

\newcommand{\rfmacrogrupo}[6]{%
  \refstepcounter{rfmacrogrupo}%
  \setcounter{rfgrupo}{0}%
  \gdef\rfmacrosigla{#2}%
  \gdef\rfmacroprioridadebase{#5}%
  \gdef\rfmacroprioridadepadrao{#5}%
  \subsection{#3 (#2)}%
  \label{#1}%
  \textbf{Escopo:} #6\\
  \textbf{Ator:} #4\\
  \textbf{Prioridade:} #5

  \vspace{0.3cm}
}

\newcommand{\rfgruporeq}[2]{%
  \refstepcounter{rfgrupo}%
  \setcounter{rfreq}{0}% -<<<<<<<<<<<<<<<<<<<<<<<<<<<<<
  \setcounter{rfcadastrotopico}{0}%
  \gdef\rfmacroprioridadepadrao{\rfmacroprioridadebase}%
  \subsubsection{#2 (\rfmacrosigla)}%
  \label{#1}%
  \vspace{0.3cm}
}

\newcommand{\rfgruporeqator}[3]{%
  \refstepcounter{rfgrupo}%
  \setcounter{rfreq}{0}%
  \gdef\rfmacroprioridadepadrao{\rfmacroprioridadebase}%
  \subsubsection{#2 (\rfmacrosigla)}%
  \label{#1}%
  \noindent\textbf{Ator:} #3

  \vspace{0.3cm}
}

\newcommand{\rfgruporeqatorprio}[4]{%
  \refstepcounter{rfgrupo}%
  \setcounter{rfreq}{0}%
  \gdef\rfmacroprioridadepadrao{#4}%
  \subsubsection{#2 (\rfmacrosigla)}%
  \label{#1}%
  \noindent\textbf{Ator:} #3\\
  \textbf{Prioridade:} #4

  \vspace{0.3cm}
}

\newcommand{\rfgruposemrequisitos}[1]{%
  \noindent\textit{#1}

  \vspace{0.3cm}
}

\newcommand{\rfcadastrotopico}[3]{%
  \refstepcounter{rfcadastrotopico}%
  \gdef\rfmacroprioridadepadrao{#3}%
  \noindent\textbf{\rfcadastrotopicoid\hspace{0.5em}#1}\par
  \noindent\textbf{Código base:} #2\\
  \textbf{Prioridade:} #3

  \vspace{0.2cm}
}

\newcommand{\rfmostraprioridade}[1]{%
  \begingroup
  \edef\rfprioridaderequisito{#1}%
  \edef\rfprioridademacro{\rfmacroprioridadepadrao}%
  \ifx\rfprioridaderequisito\rfprioridademacro
  \else
    \textbf{Prioridade específica:} #1\\
  \fi
  \endgroup
}

% #1 = rotulo de referencia (vazio para requisitos REQ-*)
% #2 = nome do requisito funcional
% #3 = descricao
% #4 = prioridade
% #5 = dependencia
% #6 = criterio de aceite
\makeatletter
\newcommand{\reqfuncfull}[6]{%
  \refstepcounter{rfreq}%
  \reqcad@split#2 \relax
  \expandafter\reqcad@testREQ\reqcad@code ZZZ\@nil
  \ifnum\reqcad@isreq=1\relax
    \edef\@currentlabel{\reqcad@code}%
    \noindent\textbf{[\reqcad@code] \reqcad@title}\\
    \label{\reqcad@code}%
    \textbf{Descrição:} #3\\
    \textbf{Dependência:} #5\\
    \textbf{Critério de aceite:} #6
  \else
    \noindent\textbf{[\rfreqid] #2}\\
    \label{#1}%
    \textbf{Descrição:} #3\\
    \rfmostraprioridade{#4}%
    \textbf{Dependência:} #5\\
    \textbf{Critério de aceite:} #6
  \fi

  \vspace{0.4cm}
}
\makeatother

\newcommand{\reqfunc}[5]{%
  \reqfuncfull{#1}{#2}{#3}{\rfmacroprioridadepadrao}{#4}{#5}%
}

\newcommand{\reqfuncprio}[6]{%
  \reqfuncfull{#1}{#2}{#3}{#4}{#5}{#6}%
}

\newcommand{\reqfunccad}[5]{%
  \reqfuncfull{}{#1}{#2}{#3}{#4}{#5}%
}

\rfmacrogrupo
{mg:autenticacao_perfil}
{AP}
{Autenticação e Perfil}
{Definido no nível do grupo}
{Essencial}
{Organizar a autenticação, a recuperação de acesso e a identificação do perfil de acesso no sistema.}
\rfgruporeqator
{g:acesso_identidade_aluno}
{Acesso e identidade do aluno}
{Aluno}

\reqfunc
{rf:autenticar_login_aluno}
{Autenticar login do aluno}
{O sistema deve autenticar o acesso do aluno.}
{Depende do requisito [\ref{REQ-CAD-ALU-01}].}
{\begin{itemize}[leftmargin=0.6cm]
    \item O sistema deve aceitar CPF ou e-mail como identificador de login do aluno.
    \item O sistema deve autenticar o aluno antes de liberar funcionalidades protegidas, impedindo acesso com credenciais inválidas, ausentes, expiradas ou não vinculadas a perfil de aluno.
    \item O sistema deve informar o motivo da falha de autenticação de forma compatível com o cenário ocorrido, sem revelar dados de outros perfis.
\end{itemize}}

\reqfunc
{rf:recuperar_senha_aluno}
{Recuperar senha do aluno}
{O sistema deve redefinir a credencial de acesso do aluno mediante verificação de identidade.}
{Depende do requisito [\ref{REQ-CAD-ALU-01}].}
{\begin{itemize}[leftmargin=0.6cm]
    \item O sistema deve aceitar e-mail como identificador para recuperação de senha do aluno.
    \item O sistema deve permitir iniciar o fluxo sem autenticação.
    \item Ao confirmar a identidade, o sistema deve exigir o CPF e a definição de nova senha antes do acesso.
    \item O sistema deve recusar identificações inválidas, inexistentes ou bloqueadas, e rejeitar validações incorretas, expiradas ou já utilizadas, sem revelar dados.
\end{itemize}}

\reqfunc
{rf:recuperar_acesso_aluno_dispositivo}
{Recuperar acesso do aluno em novo dispositivo}
{O sistema deve recuperar o acesso do aluno à conta em novo dispositivo, mediante verificação de identidade.}
{Depende do requisito [\ref{REQ-CAD-ALU-01}].}
{\begin{itemize}[leftmargin=0.6cm]
    \item O sistema deve permitir recuperação mediante confirmação de identidade.
    \item Após recuperação, o sistema deve invalidar sessões anteriores.
    \item O sistema deve impedir acesso e recusar validações inválidas, expiradas ou inconsistentes enquanto a identidade não for confirmada.
\end{itemize}}

\reqfunc
{rf:restringir_acesso_perfil_aluno}
{Restringir acesso do perfil de aluno}
{O sistema deve restringir o acesso do aluno às funcionalidades autorizadas.}
{Depende do requisito [\ref{rf:autenticar_login_aluno}].}
{\begin{itemize}[leftmargin=0.6cm]
    \item O sistema deve permitir acesso apenas às funcionalidades autorizadas ao aluno.
    \item O sistema deve negar acesso a funcionalidades fora do escopo do perfil.
    \item O sistema deve exibir mensagem de permissão insuficiente quando necessário.
    \item O sistema não deve expor dados administrativos ou de outros perfis.
\end{itemize}}

\reqfunc
{rf:direcionar_login_aluno}
{Direcionar login do aluno}
{O sistema deve direcionar o aluno autenticado diretamente ao seu dashboard, sem opção de seleção de perfil.}
{Depende do requisito [RF-AP-\ref{rf:autenticar_login_aluno}].}
{\begin{itemize}[leftmargin=0.6cm]
    \item O sistema deve redirecionar automaticamente para o dashboard do aluno após autenticação bem-sucedida.
    \item O sistema não deve apresentar opções de seleção de perfil para alunos.
\end{itemize}}
\rfgruporeqator
{g:acesso_identidade_professora}
{Acesso e identidade da professora}
{Professora}

\reqfunc
{rf:autenticar_login_professora}
{Autenticar login da professora}
{O sistema deve autenticar o acesso da professora.}
{Nenhuma.}
{\begin{itemize}[leftmargin=0.6cm]
    \item O sistema deve aceitar CPF ou e-mail como identificador de login da professora.
    \item O sistema deve autenticar a professora antes de liberar funcionalidades protegidas, impedindo acesso com credenciais inválidas, ausentes, expiradas ou não vinculadas a perfil de professora.
    \item O sistema deve exibir mensagem de rejeição compatível com o cenário de falha, sem expor dados protegidos do sistema.
\end{itemize}}

\reqfunc
{rf:restringir_acesso_perfil_professora}
{Restringir acesso do perfil de professora}
{O sistema deve restringir o acesso da professora às funcionalidades autorizadas.}
{Depende do requisito [RF-AP-\ref{rf:autenticar_login_professora}].}
{\begin{itemize}[leftmargin=0.6cm]
    \item O sistema deve restringir o perfil de professora às funcionalidades autorizadas: gestão administrativa e de repertório, manutenção de vínculos entre turma e \textit{mix}, visualização operacional nominal da turma, cancelamento administrativo de \textit{check-ins} e consulta dos painéis gerenciais.
    \item O sistema deve negar acesso a funcionalidades não declaradas para a professora, exibindo mensagem de permissão insuficiente ou de contexto não autorizado.
    \item O sistema não deve disponibilizar ações além das previstas para o perfil de professora.
\end{itemize}}

\reqfunc
{rf:direcionar_login_professora}
{Direcionar login da professora}
{O sistema deve apresentar à professora opções de acesso após autenticação, permitindo seleção entre participar como aluno, gestão de aula ou gestão administrativa.}
{Depende do requisito [RF-AP-\ref{rf:autenticar_login_professora}].}
{\begin{itemize}[leftmargin=0.6cm]
    \item O sistema deve oferecer três opções: participar como aluno (\textit{dashboard} do aluno), gestão de aula (\textit{dashboard} da professora) e gestão administrativa (\textit{dashboard} administrativo).
    \item O sistema deve permitir seleção da opção desejada pela professora.
    \item Após seleção, o sistema deve direcionar para o \textit{dashboard} correspondente.
\end{itemize}}

\rfmacrogrupo
{mg:checkin_aluno}
{CA}
{Check-in do Aluno}
{Aluno}
{Essencial}
{Organizar a jornada do aluno na consulta, reserva e gestão do próprio \textit{check-in}.}

\rfgruporeqator
{g:jornada_checkin_aluno}
{Jornada do aluno no check-in}
{Aluno}

\reqfunc
{rf:painel_aulas_dia}
{Painel geral de aulas do dia}
{O sistema deve exibir ao aluno um painel com as aulas disponíveis no dia atual.}
{Depende dos requisitos [\ref{REQ-CAD-TUR-01}] e [\ref{REQ-CAD-SAL-01}].}
{\begin{itemize}[leftmargin=0.6cm]
    \item O sistema deve disponibilizar apenas turmas da data atual.
    \item O sistema deve fornecer, para cada turma disponível, horário de início, horário de fim, professora, sala, total de vagas reserváveis e vagas ocupadas.
    \item O sistema deve desconsiderar turmas inativas, sem horário definido, sem sala válida associada ou sem professora associada.
    \item O sistema deve excluir a \textit{bike} da professora do total de vagas reserváveis.
    \item O sistema deve apresentar mensagem informativa quando não houver turmas disponíveis no dia.
\end{itemize}}

\reqfunc
{rf:detalhe_aula_selecionada}
{Exibir detalhes da aula selecionada}
{O sistema deve apresentar ao aluno os detalhes da aula selecionada.}
{Depende dos requisitos [RF-CA-\ref{rf:painel_aulas_dia}], [\ref{REQ-CAD-BIK-01}] e [\ref{REQ-CAD-MAN-01}].}
{\begin{itemize}[leftmargin=0.6cm]
    \item O sistema deve exibir horário de início, horário de fim, nome da professora e sala da turma selecionada.
    \item O sistema deve exibir os dados da turma selecionada sem misturar informações de outra turma.
    \item O sistema deve indicar a presença ou ausência de \textit{mix} vigente.
    \item Quando houver \textit{mix} vigente, o sistema deve disponibilizar acesso ao repertório e aos detalhes musicais associados à aula.
    \item O sistema deve disponibilizar o total de vagas reserváveis, vagas disponíveis e a representação das posições associadas às \textit{bikes}.
    \item O sistema deve classificar cada posição conforme seu estado operacional (`livre', `ocupada', `em manutenção', `reservada para a professora'), representar visualmente cada estado de forma distinta, e permitir seleção apenas de posições `livre'.
    \item O sistema deve identificar a posição associada ao aluno autenticado quando houver \textit{check-in} ativo.
    \item O sistema deve permitir ao aluno visualizar apenas o primeiro nome dos alunos com \textit{check-in} ativo nas posições ocupadas, sem expor sobrenome, dados de contato ou outras informações de identificação.
    \item O sistema deve excluir \textit{bikes} em manutenção e a posição da professora das vagas reserváveis e disponíveis.
    \item Se a turma deixar de estar disponível, ficar sem sala válida ou perder dados mínimos antes da abertura do painel, o sistema deve informar a indisponibilidade e retornar ao painel geral sem exibir informações desatualizadas.
\end{itemize}}

\reqfunc
{rf:realizar_checkin}
{Realizar check-in}
{O sistema deve registrar o check-in do aluno autenticado em uma vaga reservável da turma selecionada, desde que atendidas as condições de elegibilidade, disponibilidade e consistência da operação.}
{Depende dos requisitos [\ref{REQ-CAD-ALU-01}] e [RF-CA-\ref{rf:detalhe_aula_selecionada}].}
{\begin{itemize}[leftmargin=0.6cm]
    \item O sistema deve permitir o registro do \textit{check-in} somente dentro da janela de 30 minutos antes do \texttt{horarioInicio} até o próprio \texttt{horarioInicio} da turma.
    \item O sistema deve permitir reserva apenas em posição classificada como livre e reservável, recusando posições indisponíveis, incluindo aquelas associadas à professora ou em manutenção.
    \item O sistema deve impedir a realização de mais de um \textit{check-in} ativo para o mesmo aluno na mesma turma e data, bem como em turmas com sobreposição de horário.
    \item O sistema deve permitir o registro apenas quando a sessão autenticada corresponder ao aluno para o qual a reserva está sendo realizada.
    \item O sistema deve recusar a operação quando a vaga, a turma ou a janela de reserva deixarem de estar disponíveis durante o fluxo.
    \item Após registro bem-sucedido, o sistema deve atualizar imediatamente a representação das posições, a contagem de vagas e indicar a confirmação da operação ao aluno.
\end{itemize}}

\reqfunc
{rf:cancelar_checkin}
{Cancelar check-in}
{O sistema deve cancelar o check-in ativo do aluno quando solicitado por ele.}
{Depende do requisito [RF-CA-\ref{rf:realizar_checkin}].}
{\begin{itemize}[leftmargin=0.6cm]
    \item O sistema deve aceitar cancelamento apenas do \textit{check-in} ativo do próprio aluno autenticado e dentro do prazo até o \texttt{horarioInicio} da turma, recusando cancelamentos de reservas de terceiros ou fora do prazo com mensagem apropriada.
    \item O sistema deve liberar a vaga imediatamente após o cancelamento bem-sucedido.
    \item O sistema deve manter o \textit{check-in} cancelado visível com status de cancelado.
    \item Quando a reserva já não estiver ativa no momento da confirmação, o sistema deve informar o estado atualizado sem exibir mensagem de sucesso.
\end{itemize}}

\reqfunc
{rf:fila_espera}
{Entrar na fila de espera}
{O sistema deve registrar o aluno na fila de espera da turma quando ela estiver lotada. A fila serve como registro de interesse do aluno; a gestão de quem efetivamente ocupa uma vaga liberada é feita presencialmente pela professora, fora do sistema.}
{Depende dos requisitos [\ref{REQ-CAD-ALU-01}], [RF-CA-\ref{rf:detalhe_aula_selecionada}] e [RF-CA-\ref{rf:realizar_checkin}].}
{\begin{itemize}[leftmargin=0.6cm]
    \item O sistema deve aceitar no máximo 5 alunos na fila da turma, recusando novas entradas quando cheia e exibindo mensagem informativa.
    \item Se a turma deixar de estar lotada antes da conclusão da entrada, não deve registrar o aluno na fila.
    \item O sistema deve impedir duplicidade de entrada na fila para a mesma turma e data.
    \item O sistema deve impedir coexistência de fila de espera e \textit{check-in} confirmado para a mesma turma e data.
    \item O sistema deve impedir entrada da professora na fila.
    \item O sistema deve exibir a fila ordenada por posição de entrada no mapa operacional da professora, para apoio à gestão presencial.
\end{itemize}}

\reqfuncprio
{rf:consultar_repertorio_aula}
{Consultar repertório e detalhes musicais da aula}
{O sistema deve exibir ao aluno o repertório do \textit{mix} vigente e os detalhes musicais associados à aula.}
{Importante}
{Depende dos requisitos [\ref{REQ-CAD-MIX-01}], [\ref{REQ-CAD-TMX-01}] e [RF-CA-\ref{rf:detalhe_aula_selecionada}].}
{\begin{itemize}[leftmargin=0.6cm]
    \item O sistema deve disponibilizar as músicas do \textit{mix} vigente na sequência planejada.
    \item O sistema deve permitir acesso aos detalhes de cada música do repertório.
    \item Ao detalhar uma música, o sistema deve exibir os dados disponíveis de artista ou banda e categoria musical.
    \item Ao detalhar uma música, o sistema deve exibir os links de vídeo-aula associados quando houver associação válida com \textit{VideoAula}.
    \item O sistema deve sinalizar ausência de artista, categoria ou link apenas no item afetado, sem impedir a consulta das demais músicas.
    \item Quando não houver músicas no \textit{mix}, o sistema deve informar a ausência de repertório disponível.
    \item Se o \textit{mix} deixar de estar disponível ao abrir a lista, o sistema deve informar a indisponibilidade.
\end{itemize}}

\reqfuncprio
{rf:acesso_painel_participacao}
{Acesso ao painel de participação do aluno}
{O sistema deve disponibilizar ao aluno, no fluxo principal, um acesso ao painel de participação do aluno.}
{Desejável}
{Depende dos requisitos [RF-AP-\ref{rf:autenticar_login_aluno}] e [RF-AA-\ref{rf:consultar_historico_de_presenca}].}
{\begin{itemize}[leftmargin=0.6cm]
    \item O sistema deve permitir ao aluno autenticado acesso ao próprio painel de participação, encaminhando-o somente ao painel vinculado à sessão autenticada.
    \item O ponto de acesso deve exibir apenas um indicativo de entrada ao painel de acompanhamento individual.
    \item Quando não houver dados de participação, o painel de destino deve abrir com mensagem informativa.
\end{itemize}}

\reqfuncprio
{rf:avaliar_musica_repertorio}
{Registrar e modificar avaliação de música}
{O sistema deve permitir que o aluno autenticado registre ou modifique sua avaliação para uma música disponível no repertório.}
{Desejável}
{Depende dos requisitos [RF-CA-\ref{rf:consultar_repertorio_aula}] e [\ref{REQ-CAD-MUS-01}].}
{\begin{itemize}[leftmargin=0.6cm]
    \item O sistema deve permitir que o aluno registre uma avaliação para cada música disponível no repertório consultado.
    \item O sistema deve permitir a alteração posterior da avaliação atribuída pelo próprio aluno.
    \item O sistema deve manter apenas uma avaliação vigente por combinação de aluno e música.
    \item Ao consultar uma música previamente avaliada, o sistema deve apresentar a avaliação atualmente registrada pelo aluno.
    \item O sistema deve impedir que o aluno registre ou modifique avaliações em nome de outro usuário.
    \item O sistema deve informar o resultado da operação após o registro ou a atualização da avaliação.
\end{itemize}}



\rfmacrogrupo
{mg:gestao_administrativa}
{GADM}
{Gestão Administrativa}
{Professora}
{Essencial}
{Manter a base cadastral e a estrutura operacional que sustentam a operação das aulas.}

Nos subtópicos padronizados de cadastro, o item `Padronização de entradas' deve ser declarado apenas quando houver regra específica de enumeração, seleção a partir de listas controladas, cardinalidade de associação ou controle formal de entrada relevante para a entidade. Quando não houver restrição dessa natureza, o item pode ser omitido nos novos cadastros; ajustes de rastreabilidade histórica devem ser registrados em documento auxiliar, sem aparecer no documento de entrega.

\rfcadastrotopico{Manter Aluno}{REQ-CAD-ALU}{Essencial}

\reqfunccad
{REQ-CAD-ALU-01 Manter aluno}
{O sistema deve manter o cadastro de alunos, que são os participantes das turmas de spinning e cujo estado ativo determina a elegibilidade para check-in.}
{Essencial}
{Nenhuma.}
{\begin{itemize}[leftmargin=0.6cm]
    \item O sistema deve permitir cadastrar, consultar, atualizar, inativar e reativar alunos.
    \item O sistema deve exigir `nome', `cpf', `email', `dataNascimento' e `genero' como campos obrigatórios.
    \item O sistema deve permitir `telefone', `urlFoto', `instagram', `facebook', `tiktok' e `observacoes' como campos opcionais.
    \item O sistema deve validar formato de `cpf', `email' e `urlFoto' quando informada, unicidade de `cpf' e `email', consistência de `dataNascimento', e limitar `genero' às opções \texttt{MASCULINO}, \texttt{FEMININO} e \texttt{OUTRO} na inclusão e atualização.
    \item O sistema deve permitir consulta por `nome', `cpf', `email' e status, com ordenação alfabética por nome.
    \item O sistema deve preservar o histórico de participação do aluno após inativação.
    \item O sistema deve recusar a inativação quando houver tratamento operacional obrigatório pendente.
\end{itemize}}

\rfcadastrotopico{Manter Grupo de Alunos}{REQ-CAD-GAL}{Importante}

\reqfunccad
{REQ-CAD-GAL-01 Manter grupo de alunos}
{O sistema deve manter o cadastro de grupos de alunos, que organizam alunos em conjuntos para fins de gestão e comunicação operacional.}
{Importante}
{Depende do requisito [\ref{REQ-CAD-ALU-01}].}
{\begin{itemize}[leftmargin=0.6cm]
    \item O sistema deve permitir cadastrar, consultar, atualizar, inativar e reativar grupos de alunos.
    \item O sistema deve exigir `nome' como campo obrigatório nos fluxos de criação e edição.
    \item O sistema deve permitir criar grupo sem alunos associados inicialmente.
    \item O sistema deve validar unicidade de `nome' na inclusão e na atualização.
    \item O sistema deve permitir consulta por `nome' e status, com ordenação alfabética por nome.
    \item O sistema deve exibir a composição atual do grupo quando houver alunos associados.
    \item O sistema deve impedir que o mesmo aluno seja associado mais de uma vez ao mesmo grupo.
    \item O sistema deve manter o cadastro individual dos alunos inalterado ao incluir ou remover alunos do grupo.
    \item O sistema deve manter inalterado o estado cadastral dos alunos associados quando o grupo for inativado.
\end{itemize}}

\rfcadastrotopico{Manter Fabricante}{REQ-CAD-FAB}{Importante}

\reqfunccad
{REQ-CAD-FAB-01 Manter fabricante}
{O sistema deve manter o cadastro de fabricantes, que identificam a origem das bikes e são referenciados no cadastro de equipamentos.}
{Importante}
{Nenhuma.}
{\begin{itemize}[leftmargin=0.6cm]
    \item O sistema deve permitir cadastrar, consultar, atualizar, inativar e reativar fabricantes.
    \item O sistema deve exigir `nome' como campo obrigatório.
    \item O sistema deve permitir `emailContato' e `telefoneContato' como campos opcionais.
    \item O sistema deve validar a unicidade de `nome' na inclusão e na atualização.
    \item O sistema deve permitir consulta por `nome', por dados de contato informados e por status, com ordenação alfabética por nome.
\end{itemize}}

\rfcadastrotopico{Manter Tipo de Manutenção}{REQ-CAD-TMA}{Importante}

\reqfunccad
{REQ-CAD-TMA-01 Manter tipo de manutenção}
{O sistema deve manter o cadastro de tipos de manutenção, que classificam as intervenções realizadas nas bikes e determinam o registro operacional correspondente.}
{Importante}
{Nenhuma.}
{\begin{itemize}[leftmargin=0.6cm]
    \item O sistema deve permitir cadastrar, consultar, atualizar, inativar e reativar tipos de manutenção.
    \item O sistema deve exigir `nome' como campo obrigatório.
    \item O sistema deve permitir `descricao' como campo opcional.
    \item O sistema deve validar a unicidade de `nome' na inclusão e na atualização.
    \item O sistema deve permitir consulta por `nome' e por status, com ordenação alfabética por nome.
    \item O sistema deve preservar a exibição do tipo originalmente associado em manutenções já registradas, mesmo após inativação.
\end{itemize}}

\rfcadastrotopico{Manter Sala}{REQ-CAD-SAL}{Essencial}

\reqfunccad
{REQ-CAD-SAL-01 Manter sala}
{O sistema deve manter o cadastro de salas, que definem o espaço físico e a grade de posições onde as bikes são alocadas para uso nas turmas.}
{Essencial}
{Nenhuma.}
{\begin{itemize}[leftmargin=0.6cm]
    \item O sistema deve permitir cadastrar, consultar, atualizar, inativar e reativar salas.
    \item O sistema deve exigir `nome', `numeroFilas', `numeroColunas', `numeroFilaProfessora' e `numeroColunaProfessora' como campos obrigatórios nos fluxos de inclusão e atualização.
    \item O sistema deve apresentar um layout visual das posições das bikes na sala ao definir `numeroFilas' e `numeroColunas', representando cada posição como uma bike disponível.
    \item O sistema deve aceitar apenas valores válidos para a grade da sala, exigindo `numeroFilas' e `numeroColunas' positivos, e `numeroFilaProfessora' e `numeroColunaProfessora' dentro dos limites da grade física.
    \item O sistema deve definir a posição da professora (formada por `numeroFilaProfessora' e `numeroColunaProfessora') como a bike que fica inativa para check-in de alunos, pois é utilizada pela professora durante a aula.
    \item O sistema deve impedir que a grade física dependa de listas previamente cadastradas de outras entidades.
    \item O sistema deve permitir consulta por `nome' e status, com ordenação alfabética por nome.
    \item O sistema deve recusar alterações à grade da sala que invalidem posições de bike existentes ou sejam incompatíveis com turmas ativas associadas.
\end{itemize}}

\rfcadastrotopico{Manter Bike}{REQ-CAD-BIK}{Essencial}

\reqfunccad
{REQ-CAD-BIK-01 Manter bike}
{O sistema deve manter o cadastro de bikes, que representam os equipamentos físicos de spinning associados a fabricantes e posicionados nas salas para uso nas turmas.}
{Essencial}
{Depende dos requisitos [\ref{REQ-CAD-FAB-01}] e [\ref{REQ-CAD-SAL-01}].}
{\begin{itemize}[leftmargin=0.6cm]
    \item O sistema deve permitir cadastrar, consultar, atualizar, inativar e reativar bikes.
    \item O sistema deve permitir cadastrar bike sem associação física imediata à sala.
    \item O sistema deve exigir `fabricante', `nome' e `dataCadastro' como campos obrigatórios.
    \item O sistema deve validar formato e consistência cronológica de `dataCadastro'.
    \item O sistema deve validar unicidade de `nome' na inclusão e na atualização.
    \item O sistema deve permitir consulta por `nome', `fabricante' e status, com ordenação alfabética por nome.
    \item O sistema deve permitir associar ou reposicionar bike apenas em posição existente, livre e compatível com a grade da sala.
    \item O sistema deve impedir que a mesma bike ocupe mais de uma posição simultaneamente.
\end{itemize}}

\rfcadastrotopico{Manter Turma}{REQ-CAD-TUR}{Essencial}

\reqfunccad
{REQ-CAD-TUR-01 Manter turma}
{O sistema deve manter o cadastro de turmas, que definem a recorrência semanal das aulas associando sala, horário de início, duração e professora responsável.}
{Essencial}
{Depende do requisito [\ref{REQ-CAD-SAL-01}].}
{\begin{itemize}[leftmargin=0.6cm]
    \item O sistema deve permitir cadastrar, consultar, atualizar, inativar e reativar turmas.
    \item O sistema deve exigir identificação, `horarioInicio', `duracao', `diaSemana' e `sala' como campos obrigatórios nos fluxos de inclusão e atualização.
    \item O sistema deve aceitar apenas valores válidos para `duracao' (positiva) e `diaSemana' (dentro da enumeração permitida).
    \item O sistema deve registrar a turma em estado inativo quando a ativação depender de validação posterior das regras de domínio.
    \item O sistema deve permitir consulta por sala, dia da semana, horário e status, com ordenação por dia da semana e horário de início.
    \item O sistema deve recusar ativação de turma sem sala válida, sem ao menos um dia de ocorrência ou com conflito de horário em relação a outra turma ativa na mesma sala.
\end{itemize}}

\rfcadastrotopico{Manter Manutenção}{REQ-CAD-MAN}{Essencial}

\reqfunccad
{REQ-CAD-MAN-01 Manter manutenção}
{O sistema deve manter o cadastro de manutenções, que registram intervenções em bikes e afetam sua disponibilidade operacional nas salas.}
{Essencial}
{Depende dos requisitos [\ref{REQ-CAD-BIK-01}] e [\ref{REQ-CAD-TMA-01}].}
{\begin{itemize}[leftmargin=0.6cm]
    \item O sistema deve permitir cadastrar, consultar, atualizar, cancelar, inativar e reativar registros de manutenção.
    \item O sistema deve exigir `bike', `tipoManutencao', `dataSolicitacao' e `descricao' como campos obrigatórios.
    \item O sistema deve permitir registrar manutenção sem `dataRealizacao'.
    \item O sistema deve gerenciar os estados da manutenção: pendente (sem `dataRealizacao' e sem cancelamento), realizada (com `dataRealizacao' válida igual ou posterior a `dataSolicitacao' e sem cancelamento), ou cancelada (independentemente de `dataRealizacao').
    \item O sistema deve gerar automaticamente código único e inalterável para cada novo registro de manutenção.
    \item O sistema deve permitir selecionar `bike' e `tipoManutencao' apenas dentre registros válidos e ativos apresentados pelo sistema.
    \item O sistema deve permitir consulta por `bike', `tipoManutencao', período de solicitação e estado operacional do registro.
    \item O sistema deve permitir filtros pelos estados \textbf{pendente}, \textbf{realizada} e \textbf{cancelada}.
    \item O sistema deve ordenar consultas por `dataSolicitacao' em ordem cronológica decrescente.
    \item O sistema deve refletir automaticamente no mapa da sala e na disponibilidade da turma os efeitos de manutenção pendente, realizada, cancelada ou associada a bike inativa.
    \item O sistema deve impedir que bikes com manutenção pendente apareçam como disponíveis.
\end{itemize}}

\rfmacrogrupo
{mg:gestao_aula}
{GDA}
{Gestão de aula}
{Professora}
{Essencial}
{Organizar a operação da aula, incluindo visualização e cancelamento administrativo de \textit{check-ins}, painéis de ocupação e frequência, e manter o repertório e os vínculos que definem seu uso por turma.}

\rfgruporeq
{g:operacao_aula}
{Operação da aula e painéis gerenciais}

\reqfuncfull
{rf:visualizar_mapa_operacional}
{Visualizar mapa operacional da turma}
{O sistema deve exibir à professora o mapa operacional da turma selecionada.}
{Importante}
{Depende dos requisitos [RF-AP-\ref{rf:restringir_acesso_perfil_professora}], [\ref{REQ-CAD-TUR-01}], [RF-CA-\ref{rf:realizar_checkin}] e [RF-CA-\ref{rf:fila_espera}].}
{\begin{itemize}[leftmargin=0.6cm]
    \item O sistema deve disponibilizar a representação da ocupação efetiva da turma na data consultada.
    \item O sistema deve associar os alunos às posições ocupadas, distinguindo os estados: confirmado, em fila de espera e cancelado.
    \item Quando não houver reservas ativas, o sistema deve exibir a configuração da sala sem nomes.
    \item Quando não houver sala válida ou contexto montável, o sistema deve informar a indisponibilidade.
    \item O sistema deve exibir discretamente a idade média dos alunos com \textit{check-in} ativo na turma; quando não houver nenhum \textit{check-in} ativo, o indicador não deve ser exibido.
\end{itemize}}

\reqfuncfull
{rf:cancelar_checkins_de_alunos}
{Cancelar \textit{check-ins} de alunos}
{O sistema deve cancelar check-ins de alunos da turma no contexto operacional, mediante ação da professora.}
{Importante}
{Depende dos requisitos [RF-AP-\ref{rf:restringir_acesso_perfil_professora}], [RF-GDA-\ref{rf:visualizar_mapa_operacional}] e [RF-CA-\ref{rf:cancelar_checkin}].}
{\begin{itemize}[leftmargin=0.6cm]
    \item Mediante ação da professora, o sistema deve permitir seleção apenas de posições ocupadas com \textit{check-in} ativo para cancelamento.
    \item O sistema deve impedir seleção de posições livres, em manutenção ou reservadas para a professora.
    \item Após o cancelamento, o sistema deve liberar a vaga imediatamente.
    \item O sistema deve manter o \textit{check-in} cancelado visível com o estado de cancelado.
    \item Se a reserva não estiver mais ativa no momento da confirmação, o sistema deve informar o estado atual e atualizar o mapa sem exibir sucesso indevido.
    \item O sistema deve recusar tentativas de cancelamento fora do contexto da turma ou do período permitido, exibindo mensagem compatível.
\end{itemize}}

\reqfuncfull
{rf:consultar_historico_de_mixes}
{Consultar histórico de mixes da turma}
{O sistema deve exibir à professora o histórico de \textit{mixes} utilizados pela turma.}
{Desejável}
{Depende dos requisitos [RF-AP-\ref{rf:restringir_acesso_perfil_professora}] e [\ref{REQ-CAD-TMX-01}].}
{\begin{itemize}[leftmargin=0.6cm]
    \item O sistema deve listar os \textit{mixes} encerrados da turma em ordem cronológica decrescente, exibindo nome e intervalo de uso de cada vínculo encerrado.
    \item Quando não houver histórico encerrado, o sistema deve informar a ausência e manter acessível o \textit{mix} vigente, quando houver.
    \item Quando a turma nunca tiver tido \textit{mix} associado, o sistema deve informar a ausência sem apresentar dados de outras turmas.
    \item O histórico deve refletir a vigência do vínculo `TurmaMix', não a vigência geral de `Mix'.
\end{itemize}}

\reqfuncfull
{rf:exibir_painel_ocupacao}
{Exibir painel de ocupação do dia}
{O sistema deve exibir à professora um painel de ocupação das turmas do dia atual.}
{Importante}
{Depende dos requisitos [RF-AP-\ref{rf:restringir_acesso_perfil_professora}], [\ref{REQ-CAD-TUR-01}], [RF-CA-\ref{rf:realizar_checkin}] e [RF-CA-\ref{rf:cancelar_checkin}].}
{\begin{itemize}[leftmargin=0.6cm]
    \item O sistema deve exibir todas as turmas ativas do dia com a respectiva taxa de ocupação, incluindo turmas sem \textit{check-ins} com ocupação zero.
    \item Quando não houver turmas ativas no dia, o sistema deve exibir mensagem informativa sem apresentar dados parciais.
    \item Ao selecionar uma turma, o sistema deve exibir o mapa operacional correspondente.
    \item Se a turma deixar de estar disponível entre a carga do painel e a seleção, o sistema deve informar a indisponibilidade e atualizar a listagem sem exibir dados desatualizados.
    \item Quando o sistema carregar o painel mas não conseguir calcular a ocupação de uma turma específica, o sistema deve exibir essa turma com indicação de dado indisponível e manter as demais turmas visíveis com seus dados calculados.
\end{itemize}}

\reqfuncfull
{rf:exibir_painel_frequencia}
{Exibir painel de frequência dos alunos}
{O sistema deve exibir à professora um painel de frequência individual dos alunos no período selecionado.}
{Importante}
{Depende dos requisitos [RF-AP-\ref{rf:restringir_acesso_perfil_professora}], [RF-CA-\ref{rf:realizar_checkin}] e [RF-CA-\ref{rf:cancelar_checkin}].}
{\begin{itemize}[leftmargin=0.6cm]
    \item O sistema deve suportar filtro por aluno, turma e período.
    \item A contagem de frequência deve incluir apenas \textit{check-ins} confirmados, excluindo cancelados.
    \item Cada linha do painel deve exibir o nome do aluno, o total de \textit{check-ins} confirmados e o percentual calculado exclusivamente sobre as aulas da turma filtrada no período selecionado.
    \item O sistema deve permitir ordenação dos resultados por nome ou por frequência.
    \item Quando não houver dados no período selecionado, o sistema deve exibir mensagem informativa sem apresentar linhas com valores zerados ou vazios.
    \item Quando o sistema carregar o painel mas não conseguir calcular a frequência para uma combinação específica de filtros, o sistema deve sinalizar o dado indisponível naquela linha sem ocultar os demais resultados calculados.
\end{itemize}}

\reqfuncfull
{rf:consultar_avaliacoes_musicas}
{Consultar avaliações das músicas}
{O sistema deve permitir que a professora consulte as avaliações registradas pelos alunos para as músicas do repertório.}
{Importante}
{Depende dos requisitos [RF-CA-\ref{rf:avaliar_musica_repertorio}], [\ref{REQ-CAD-MUS-01}] e [RF-AP-\ref{rf:restringir_acesso_perfil_professora}].}
{\begin{itemize}[leftmargin=0.6cm]
    \item O sistema deve permitir a consulta das avaliações associadas às músicas do repertório.
    \item O sistema deve apresentar, para cada música, a quantidade de avaliações registradas e a avaliação média correspondente.
    \item O sistema deve permitir consulta por música, artista ou banda e mix.
    \item O sistema deve permitir ordenar os resultados pela avaliação média ou pela quantidade de avaliações registradas.
    \item Quando uma música não possuir avaliações, o sistema deve apresentar essa informação sem atribuir nota média artificial.
    \item Quando não houver resultados compatíveis com os filtros aplicados, o sistema deve exibir mensagem informativa.
\end{itemize}}

\reqfuncfull
{rf:copiar_marcacoes_instagram_alunos_turma}
{Copiar marcações de alunos para publicação no Instagram}
{O sistema deve permitir que a professora copie a relação de perfis do Instagram dos alunos vinculados a uma turma, para utilização em publicações realizadas fora do sistema.}
{Desejável}
{Depende dos requisitos [\ref{REQ-CAD-TUR-01}], [\ref{REQ-CAD-ALU-01}] e [RF-AP-\ref{rf:restringir_acesso_perfil_professora}].}
{\begin{itemize}[leftmargin=0.6cm]
    \item O sistema deve disponibilizar a funcionalidade a partir da consulta de uma turma.
    \item O sistema deve apresentar uma opção específica para visualizar a relação de perfis do Instagram dos alunos vinculados à turma selecionada.
    \item O sistema deve incluir somente alunos que possuam perfil do Instagram informado em seu cadastro.
    \item O sistema deve apresentar os perfis em formato textual adequado para marcação em publicações, utilizando o prefixo \texttt{@}.
    \item O sistema deve permitir copiar integralmente a relação de perfis para a área de transferência do dispositivo.
    \item O sistema deve informar à professora quando a cópia for concluída com sucesso.
    \item Quando nenhum aluno da turma possuir perfil do Instagram informado, o sistema deve apresentar mensagem informativa.
\end{itemize}}

\rfgruporeq
{g:cadastros_repertorio}
{Cadastros musicais do repertório}

\rfcadastrotopico{Manter Artista ou Banda}{REQ-CAD-ARB}{Importante}

\reqfunccad
{REQ-CAD-ARB-01 Manter artista ou banda}
{O sistema deve manter o cadastro de artistas e bandas, referenciados nas músicas do repertório como responsáveis pelas obras.}
{Importante}
{Nenhuma.}
{\begin{itemize}[leftmargin=0.6cm]
    \item O sistema deve permitir cadastrar, consultar, atualizar, inativar e reativar artistas e bandas.
    \item O sistema deve exigir `nome' como campo obrigatório nos fluxos de inclusão e atualização.
    \item O sistema deve permitir cadastro sem `descricao', `link' ou `urlFoto'.
    \item O sistema deve validar formato de `link' e `urlFoto' quando informados, e unicidade de `nome' na inclusão e atualização.
    \item O sistema deve permitir consulta por `nome' e status, com ordenação alfabética por nome.
\end{itemize}}

\rfcadastrotopico{Manter Categoria Musical}{REQ-CAD-CAT}{Importante}

\reqfunccad
{REQ-CAD-CAT-01 Manter categoria musical}
{O sistema deve manter categorias musicais e suas associações às músicas cadastradas, permitindo classificar o repertório por estilo ou gênero sem sobrecarregar o cadastro principal de música.}
{Importante}
{Depende do requisito [\ref{REQ-CAD-MUS-01}].}
{\begin{itemize}[leftmargin=0.6cm]
    \item O sistema deve permitir cadastrar, consultar, atualizar, inativar e reativar categorias musicais.
    \item O sistema deve permitir buscar uma música cadastrada por digitação, apresentando opções correspondentes para seleção controlada.
    \item O sistema deve exigir `nome' como campo obrigatório nos fluxos de inclusão e atualização.
    \item O sistema deve permitir cadastro sem `descricao'.
    \item O sistema deve validar unicidade de `nome' na inclusão e na atualização.
    \item O sistema deve permitir selecionar múltiplas categorias para a música escolhida, adicionando cada categoria selecionada a um painel de resultados da seleção.
    \item Quando a categoria digitada não existir, o sistema deve permitir cadastrá-la como nova categoria antes de associá-la à música.
    \item O sistema deve permitir remover categorias do painel de seleção antes da confirmação.
    \item O sistema não deve permitir associar a mesma categoria musical mais de uma vez à mesma música.
    \item O sistema deve permitir consulta por `nome' e status, com ordenação alfabética por nome.
\end{itemize}}

\rfcadastrotopico{Manter Vídeo de Aula}{REQ-CAD-VID}{Importante}

\reqfunccad
{REQ-CAD-VID-01 Manter vídeo de aula}
{O sistema deve manter links de vídeos de aula e suas associações às músicas cadastradas, para consulta pelos alunos durante as aulas.}
{Importante}
{Depende do requisito [\ref{REQ-CAD-MUS-01}].}
{\begin{itemize}[leftmargin=0.6cm]
    \item O sistema deve permitir cadastrar, consultar, atualizar, inativar e reativar vídeos de aula.
    \item O sistema deve permitir buscar uma música cadastrada por digitação, apresentando opções correspondentes para seleção controlada.
    \item O sistema deve exigir `nome' e `linkVideo' como campos obrigatórios nos fluxos de inclusão e atualização, sendo `linkVideo' uma URL para vídeo externo, não envio de arquivo.
    \item O sistema deve validar formato de `linkVideo'.
    \item O sistema deve validar unicidade de `linkVideo' na inclusão e na atualização.
    \item O sistema deve permitir selecionar múltiplos links de vídeo para a música escolhida, adicionando cada link selecionado a um painel de resultados da seleção.
    \item Quando a URL digitada não existir, o sistema deve permitir cadastrá-la como novo vídeo de aula antes de associá-la à música.
    \item O sistema deve permitir remover vídeos de aula do painel de seleção antes da confirmação.
    \item O sistema não deve permitir associar o mesmo vídeo de aula mais de uma vez à mesma música.
    \item O sistema deve permitir consulta por `nome', `linkVideo' e status, com ordenação alfabética por nome.
\end{itemize}}

\rfcadastrotopico{Manter Música}{REQ-CAD-MUS}{Essencial}

\reqfunccad
{REQ-CAD-MUS-01 Manter música}
{O sistema deve manter o cadastro de músicas, que compõem o repertório dos mixes e são exibidas aos alunos durante a consulta da aula.}
{Essencial}
{Depende do requisito [\ref{REQ-CAD-ARB-01}].}
{\begin{itemize}[leftmargin=0.6cm]
    \item O sistema deve permitir cadastrar, consultar, atualizar, inativar e reativar músicas.
    \item O sistema deve exigir `nome' e `artistaBanda' como campos obrigatórios na inclusão e na atualização.
    \item O sistema deve permitir cadastrar música sem categorias musicais e sem vídeos de aula associados, pois essas associações são mantidas em fluxos próprios.
    \item O sistema deve validar a unicidade da combinação de `nome' e `artistaBanda' na inclusão e na atualização.
    \item O sistema deve permitir selecionar `artistaBanda' por seleção única controlada.
    \item O sistema deve permitir consulta por `nome', `artistaBanda', `categoriaMusical' e status, com ordenação alfabética por nome.
    \item O sistema deve manter mixes já cadastrados exibindo a referência da música após atualizações válidas ou inativação.
\end{itemize}}

\rfcadastrotopico{Manter Mix}{REQ-CAD-MIX}{Essencial}

\reqfunccad
{REQ-CAD-MIX-01 Manter mix}
{O sistema deve manter o cadastro de mixes, que são coleções ordenadas de músicas vinculadas a turmas por períodos de vigência para uso nas aulas.}
{Essencial}
{Depende do requisito [\ref{REQ-CAD-MUS-01}].}
{\begin{itemize}[leftmargin=0.6cm]
    \item O sistema deve permitir cadastrar, consultar, atualizar, inativar e reativar mixes.
    \item O sistema deve exigir `nome', `descricao' e ao menos uma música com ordem definida como campos obrigatórios.
    \item O sistema deve exigir ordem completa e inequívoca entre as músicas associadas na inclusão e na atualização.
    \item O sistema deve permitir cadastrar mix sem `dataFim'.
    \item O sistema deve validar unicidade de `nome', e que `dataFim' seja igual ou posterior a `dataInicio' quando informada, na inclusão e atualização.
    \item O sistema deve permitir selecionar músicas e definir a ordem de execução por controles da interface.
    \item O sistema deve permitir consulta por `nome', período de vigência e status, com ordenação cronológica decrescente por `dataInicio'.
\end{itemize}}

\rfgruporeq
{g:vinculos_vigencia_aula}
{Vínculo entre turma e mix e vigência do repertório}

\rfcadastrotopico{Manter Vínculo TurmaMix}{REQ-CAD-TMX}{Essencial}

\reqfunccad
{REQ-CAD-TMX-01 Manter vínculo TurmaMix}
{O sistema deve manter os vínculos entre turma e mix, que determinam qual repertório está ativo em cada turma em cada período de vigência.}
{Essencial}
{Depende dos requisitos [\ref{REQ-CAD-TUR-01}] e [\ref{REQ-CAD-MIX-01}].}
{\begin{itemize}[leftmargin=0.6cm]
    \item O sistema deve permitir cadastrar, consultar, atualizar, encerrar, inativar e reativar vínculos TurmaMix.
    \item O sistema deve exigir `turma', `mix' e `dataInicio' como campos obrigatórios.
    \item O sistema deve permitir cadastrar vínculo sem `dataFim'.
    \item O sistema deve exigir que `dataFim' seja igual ou posterior a `dataInicio' quando informada.
    \item O sistema não deve permitir que a mesma turma possua mais de um vínculo ativo de mix simultaneamente.
    \item O sistema deve permitir consulta por turma, mix, período e status, com ordenação cronológica decrescente por `dataInicio'.
    \item O sistema deve recuperar vínculos ativos, inativos e encerrados nas consultas.
    \item O sistema deve permitir identificar com clareza o mix vigente da turma e seus períodos anteriores, sem ambiguidade.
\end{itemize}}

\rfmacrogrupo
{mg:acompanhamento_aluno}
{AA}
{Acompanhamento do Aluno}
{Aluno}
{Importante}
{Nesta versão, o acompanhamento individual do aluno cobre a consulta ao próprio histórico de \textit{check-ins}, com acesso restrito ao aluno autenticado.}

\rfgruporeqatorprio
{g:historico_metricas_aluno}
{Histórico individual de participação}
{Aluno}
{Desejável}

\reqfuncprio
{rf:consultar_historico_de_presenca}
{Consultar histórico de presença}
{O sistema deve exibir ao aluno autenticado o próprio histórico de \textit{check-ins}.}
{Importante}
{Depende dos requisitos [RF-CA-\ref{rf:realizar_checkin}] e [RF-CA-\ref{rf:cancelar_checkin}].}
{\begin{itemize}[leftmargin=0.6cm]
    \item O sistema deve reunir apenas os \textit{check-ins} do aluno autenticado na consulta.
    \item O sistema deve distinguir \textit{check-ins} futuros de concluídos e exibir cancelados com o respectivo estado.
    \item Quando não houver histórico, o sistema deve exibir mensagem informativa.
    \item O sistema deve recusar acesso ao histórico sem autenticação válida.
    \item O sistema não deve exibir o histórico de \textit{check-ins} de terceiros.
\end{itemize}}

\subsection{Matriz de Dependências dos Requisitos Funcionais}

A Tabela~\ref{tab:matriz_dependencia} consolida as dependências declaradas entre os requisitos funcionais, após a verificação de consistência entre os cadastros, os fluxos de interação do aluno, as operações realizadas pela professora e as funcionalidades de acompanhamento individual.

Cada requisito também foi registrado como uma \textit{Issue} no GitHub. O acesso à respectiva \textit{Issue} pode ser realizado por meio do identificador apresentado na coluna \textit{Issue} da tabela, substituindo-se \texttt{XX} pelo número correspondente no seguinte endereço:

\url{https://github.com/heliokamakawa/spin_flow_pt_br/issues/XX}

\begingroup
\footnotesize
\begin{longtable}{|c|p{6cm}|p{7cm}|} 
    \caption{Matriz de Dependências dos Requisitos Funcionais.}
    \label{tab:matriz_dependencia} \\ \hline
    \textbf{Issue} & \textbf{Requisito} & \textbf{Dependência de requisitos} \\ \hline
    \endfirsthead 
    
    \caption[]{Matriz de Dependências dos Requisitos Funcionais (continuação).} \\ \hline
    \textbf{Issue} & \textbf{Requisito} & \textbf{Dependência de requisitos} \\ \hline
    \endhead

    1 & [RF-AP-\ref{rf:autenticar_login_aluno}] Autenticar login do aluno & [\ref{REQ-CAD-ALU-01}]. \\ \hline
    2 & [RF-AP-\ref{rf:recuperar_senha_aluno}] Recuperar senha do aluno & [\ref{REQ-CAD-ALU-01}]. \\ \hline
    3 & [RF-AP-\ref{rf:recuperar_acesso_aluno_dispositivo}] Recuperar acesso do aluno em novo dispositivo & [\ref{REQ-CAD-ALU-01}]. \\ \hline
    4 & [RF-AP-\ref{rf:restringir_acesso_perfil_aluno}] Restringir acesso do perfil de aluno & [RF-AP-\ref{rf:autenticar_login_aluno}]. \\ \hline
    5 & [RF-AP-\ref{rf:direcionar_login_aluno}] Direcionar login do aluno & [RF-AP-\ref{rf:autenticar_login_aluno}]. \\ \hline
    6 & [RF-AP-\ref{rf:autenticar_login_professora}] Autenticar login da professora & Nenhuma. \\ \hline
    7 & [RF-AP-\ref{rf:restringir_acesso_perfil_professora}] Restringir acesso do perfil de professora & [RF-AP-\ref{rf:autenticar_login_professora}]. \\ \hline
    8 & [RF-AP-\ref{rf:direcionar_login_professora}] Direcionar login da professora & [RF-AP-\ref{rf:autenticar_login_professora}]. \\ \hline
    9 & [RF-CA-\ref{rf:painel_aulas_dia}] Painel geral de aulas do dia & [\ref{REQ-CAD-TUR-01}] e [\ref{REQ-CAD-SAL-01}]. \\ \hline
    10 & [RF-CA-\ref{rf:detalhe_aula_selecionada}] Exibir detalhes da aula selecionada & [RF-CA-\ref{rf:painel_aulas_dia}], [\ref{REQ-CAD-BIK-01}] e [\ref{REQ-CAD-MAN-01}]. \\ \hline
    11 & [RF-CA-\ref{rf:realizar_checkin}] Realizar \textit{check-in} & [\ref{REQ-CAD-ALU-01}] e [RF-CA-\ref{rf:detalhe_aula_selecionada}]. \\ \hline
    12 & [RF-CA-\ref{rf:cancelar_checkin}] Cancelar \textit{check-in} & [RF-CA-\ref{rf:realizar_checkin}]. \\ \hline
    13 & [RF-CA-\ref{rf:fila_espera}] Entrar na fila de espera & [\ref{REQ-CAD-ALU-01}], [RF-CA-\ref{rf:detalhe_aula_selecionada}] e [RF-CA-\ref{rf:realizar_checkin}]. \\ \hline
    14 & [RF-CA-\ref{rf:consultar_repertorio_aula}] Consultar repertório e detalhes musicais da aula & [\ref{REQ-CAD-MIX-01}], \ref{REQ-CAD-TMX-01}] e [RF-CA-\ref{rf:detalhe_aula_selecionada}]. \\ \hline
    15 & [RF-CA-\ref{rf:acesso_painel_participacao}] Acesso ao painel de participação do aluno & [RF-AP-\ref{rf:autenticar_login_aluno}] e [RF-AA-\ref{rf:consultar_historico_de_presenca}]. \\ \hline
    44 & [RF-CA-\ref{rf:avaliar_musica_repertorio}] Registrar e modificar avaliação de música & [RF-CA-\ref{rf:consultar_repertorio_aula}] e [\ref{REQ-CAD-MUS-01}]. \\ \hline
    16 & [\ref{REQ-CAD-ALU-01}] Manter aluno & Nenhuma. \\ \hline
    17 & [\ref{REQ-CAD-GAL-01}] Manter grupo de alunos & [\ref{REQ-CAD-ALU-01}]. \\ \hline
    18 & [\ref{REQ-CAD-FAB-01}] Manter fabricante & Nenhuma. \\ \hline
    19 & [\ref{REQ-CAD-TMA-01}] Manter tipo de manutenção & Nenhuma. \\ \hline
    20 & [\ref{REQ-CAD-SAL-01}] Manter sala & Nenhuma. \\ \hline
    21 & [\ref{REQ-CAD-BIK-01}] Manter \textit{bike} & [\ref{REQ-CAD-FAB-01}] e [\ref{REQ-CAD-SAL-01}]. \\ \hline
    22 & [\ref{REQ-CAD-TUR-01}] Manter turma & [\ref{REQ-CAD-SAL-01}]. \\ \hline
    23 & [\ref{REQ-CAD-MAN-01}] Manter manutenção & [\ref{REQ-CAD-BIK-01}] e [\ref{REQ-CAD-TMA-01}]. \\ \hline
    24 & [RF-GDA-\ref{rf:visualizar_mapa_operacional}] Visualizar mapa operacional da turma & [RF-AP-\ref{rf:restringir_acesso_perfil_professora}], [\ref{REQ-CAD-TUR-01}], [RF-CA-\ref{rf:realizar_checkin}] e [RF-CA-\ref{rf:fila_espera}]. \\ \hline
    25 & [RF-GDA-\ref{rf:cancelar_checkins_de_alunos}] Cancelar \textit{check-ins} de alunos & [RF-AP-\ref{rf:restringir_acesso_perfil_professora}], [RF-GDA-\ref{rf:visualizar_mapa_operacional}] e [RF-CA-\ref{rf:cancelar_checkin}]. \\ \hline
    26 & [RF-GDA-\ref{rf:consultar_historico_de_mixes}] Consultar histórico de mixes da turma & [RF-AP-\ref{rf:restringir_acesso_perfil_professora}] e [\ref{REQ-CAD-TMX-01}]. \\ \hline
    27 & [RF-GDA-\ref{rf:exibir_painel_ocupacao}] Exibir painel de ocupação do dia & [RF-AP-\ref{rf:restringir_acesso_perfil_professora}], [\ref{REQ-CAD-TUR-01}], [RF-CA-\ref{rf:realizar_checkin}] e [RF-CA-\ref{rf:cancelar_checkin}]. \\ \hline
    28 & [RF-GDA-\ref{rf:exibir_painel_frequencia}] Exibir painel de frequência dos alunos & [RF-AP-\ref{rf:restringir_acesso_perfil_professora}], [RF-CA-\ref{rf:realizar_checkin}] e [RF-CA-\ref{rf:cancelar_checkin}]. \\ \hline
    45 & [RF-GDA-\ref{rf:consultar_avaliacoes_musicas}] Consultar avaliações das músicas & [RF-CA-\ref{rf:avaliar_musica_repertorio}], [\ref{REQ-CAD-MUS-01}] e [RF-AP-\ref{rf:restringir_acesso_perfil_professora}]. \\ \hline
    46 & [RF-GDA-\ref{rf:copiar_marcacoes_instagram_alunos_turma}] Copiar marcações de alunos para publicação no Instagram &  [\ref{REQ-CAD-TUR-01}], [\ref{REQ-CAD-ALU-01}] e [RF-AP-\ref{rf:restringir_acesso_perfil_professora}]. \\ \hline
    29 & [\ref{REQ-CAD-ARB-01}] Manter artista ou banda & Nenhuma. \\ \hline
    30 & [\ref{REQ-CAD-CAT-01}] Manter categoria musical & [\ref{REQ-CAD-MUS-01}]. \\ \hline
    31 & [\ref{REQ-CAD-VID-01}] Manter vídeo de aula & [\ref{REQ-CAD-MUS-01}]. \\ \hline
    32 & [\ref{REQ-CAD-MUS-01}] Manter música & [\ref{REQ-CAD-ARB-01}]. \\ \hline
    33 & [\ref{REQ-CAD-MIX-01}] Manter mix & [\ref{REQ-CAD-MUS-01}]. \\ \hline
    34 & [\ref{REQ-CAD-TMX-01}] Manter vínculo TurmaMix & [\ref{REQ-CAD-TUR-01}] e [\ref{REQ-CAD-MIX-01}]. \\ \hline
    35 & [RF-AA-\ref{rf:consultar_historico_de_presenca}] Consultar histórico de presença & [RF-CA-\ref{rf:realizar_checkin}] e [RF-CA-\ref{rf:cancelar_checkin}]. \\ \hline
\end{longtable}
\endgroup

% \input{sections/06-requisitos_funcionais_validar}

% --------------------------------------------------
% 5 REQUISITOS NÃO-FUNCIONAIS
% --------------------------------------------------
\section{Requisitos Não Funcionais}

\subsubsection*{Desempenho e Fluidez}
\reqfuncfull
{nf:concluir_reserva_em_ate_dois_segundos}
{Concluir reserva por check-in em até dois segundos}
{A operação completa de reserva por check-in, desde a validação da posição escolhida até a atualização do mapa e da disponibilidade da turma, deve concluir em no máximo dois segundos. Esse requisito se aplica ao fluxo crítico do aluno, inclusive em situações de sala com ocupação total e tentativas concorrentes de reserva no mesmo horário. O objetivo é preservar resposta rápida e previsível no principal fluxo de valor do sistema.}
{Importante}
{Depende dos requisitos funcionais do fluxo de reserva por check-in e da execução consistente das validações e atualizações correlatas. Depende de processamento e persistência capazes de atender ao tempo máximo definido mesmo em cenário de alta ocupação.}
{Contexto principal: fluxo do aluno. O requisito é mensurável e deve ser verificado com critério objetivo de tempo fim a fim da operação.}
\reqfuncfull
{nf:renderizar_mapa_em_ate_um_segundo}
{Renderizar mapa da sala em até um segundo}
{O mapa visual da sala deve ser renderizado em no máximo um segundo após a seleção da turma e da data, para salas com até cinquenta posições. Esse requisito garante que a consulta espacial da aula permaneça fluida e útil no momento de escolha da posição. O tempo de resposta deve considerar a representação da grade e dos estados das posições no contexto consultado.}
{Importante}
{Depende dos requisitos funcionais de seleção de turma e data, exibição do mapa e identificação de estados das posições. Depende da capacidade de montagem do mapa dentro do limite de desempenho definido para o tamanho de sala previsto.}
{Contexto principal: mapa de bikes. O requisito é mensurável e deve ser validado em salas com até cinquenta posições.}
\reqfuncfull
{nf:carregar_listagens_com_fluidez}
{Carregar listagens e consultas com fluidez}
{As telas de listagem e consulta do produto devem carregar e renderizar até quinhentos registros em no máximo dois segundos, adotando paginação ou carregamento progressivo quando necessário para preservar fluidez de uso. O requisito se aplica a consultas operacionais e informacionais previstas no recorte do sistema. A estratégia de exibição pode variar, desde que o comportamento percebido pelo usuário continue rápido e controlado.}
{Importante}
{Depende dos requisitos de consulta e listagem do sistema e da estratégia de exibição adotada para volumes maiores de dados. Depende de mecanismos de paginação ou carregamento progressivo quando o volume exigir.}
{Contexto principal: consultas e listagens. O requisito é mensurável e deve considerar tanto carregamento quanto renderização perceptível na interface.}

\subsubsection*{Consistência Operacional e Acessibilidade Visual}
\reqfuncfull
{nf:garantir_consistencia_da_reserva}
{Garantir consistência entre reserva, ocupação e disponibilidade}
{O sistema deve garantir consistência entre reserva, ocupação da posição e disponibilidade da turma, evitando estados contraditórios após efetivação, cancelamento ou concorrência no check-in. Toda alteração operacional relevante deve produzir resultado coerente entre o registro persistido, o mapa exibido e a vaga considerada disponível. O requisito protege a integridade do fluxo principal da aula.}
{Essencial}
{Depende dos requisitos funcionais de reserva, cancelamento, atualização de mapa e reflexo operacional de manutenção. Depende da execução coordenada das escritas e recálculos relacionados ao estado da aula.}
{Contexto principal: operação de reserva e disponibilidade. O requisito é transversal aos fluxos que alteram ocupação ou vagas da turma.}
\reqfuncfull
{nf:tratar_concorrencia_de_reserva}
{Tratar concorrência na reserva da mesma posição}
{Quando houver tentativa simultânea de reserva da mesma posição na mesma turma e data, apenas uma operação deve ser bem-sucedida e a outra deve receber retorno claro de indisponibilidade. O sistema deve tratar a concorrência como situação controlada, sem gerar dupla ocupação da posição nem inconsistências entre reserva e mapa. O requisito protege a unicidade operacional da vaga em cenários de disputa simultânea.}
{Essencial}
{Depende dos requisitos funcionais de impedimento de dupla ocupação da posição, exibição de erro na reserva e consistência entre reserva e disponibilidade. Depende da execução atômica e sincronizada da confirmação da reserva.}
{Contexto principal: concorrência de check-in. O comportamento deve ser verificável em tentativas simultâneas sobre a mesma posição.}
\reqfuncfull
{nf:nao_depender_apenas_de_cor}
{Não depender apenas de cor para estados críticos}
{O produto não deve transmitir estados críticos apenas por cor. Disponibilidade, bloqueio, ocupação e indisponibilidade devem possuir reforço adicional por texto, ícone ou padrão visual equivalente, de modo que a compreensão do estado não dependa exclusivamente da percepção cromática do usuário. O requisito se aplica especialmente ao mapa da sala e aos pontos de feedback operacional da interface.}
{Importante}
{Depende dos requisitos funcionais de identificação de estados do mapa e visualização operacional da sala. Depende da definição de elementos visuais complementares para representar os estados críticos.}
{Contexto principal: interface e acessibilidade visual. O requisito é transversal a telas com estados operacionais relevantes.}

\subsubsection*{Privacidade e Proteção de Dados}
\reqfuncfull
{nf:tratar_dados_com_finalidade_explicita}
{Tratar dados pessoais com finalidade explícita}
{Dados pessoais dos alunos devem ser tratados com finalidade explícita e não devem ser compartilhados com terceiros por meio do sistema sem base apropriada. O uso dessas informações deve permanecer vinculado ao contexto operacional e informacional previsto pelo produto. O requisito estabelece limite de finalidade e compartilhamento no tratamento dos dados pessoais.}
{Essencial}
{Depende dos requisitos funcionais que lidam com dados pessoais, históricos e identificação de alunos. Depende das regras de domínio e conformidade aplicáveis ao uso dessas informações.}
{Contexto principal: privacidade e governança de dados. O requisito é transversal a todo fluxo que trate dados pessoais do aluno.}
\reqfuncfull
{nf:restringir_visualizacao_de_dados_pessoais}
{Restringir visualização de dados pessoais e históricos}
{Informações pessoais e históricas do aluno devem ser visíveis apenas ao próprio usuário e aos perfis autorizados pelas regras do domínio. A interface e os serviços de consulta devem respeitar essa limitação de visibilidade em histórico, mapa e demais pontos de acesso a dados sensíveis. O requisito protege privacidade e controle de acesso no recorte funcional do sistema.}
{Essencial}
{Depende dos requisitos funcionais de histórico individual, anonimização de ocupações de terceiros e visualização operacional pela professora. Depende das regras de visibilidade por perfil definidas no domínio.}
{Contexto principal: controle de acesso a informações pessoais. O requisito admite exceções apenas para perfis autorizados pelo domínio.}
\reqfuncfull
{nf:armazenar_dados_em_area_protegida}
{Armazenar dados em área protegida do dispositivo}
{O armazenamento de dados do aplicativo deve ocorrer em área protegida do ambiente de execução, sem exposição deliberada de dados pessoais em locais públicos do dispositivo. Esse requisito reforça a proteção física e lógica das informações mantidas localmente, especialmente no contexto de operação prioritária offline. O objetivo é reduzir risco de acesso indevido por armazenamento inseguro.}
{Importante}
{Depende do requisito [\ref{nf:operar_prioritariamente_offline}] e do tratamento de dados pessoais previsto no sistema. Depende da escolha de mecanismos de armazenamento compatíveis com área protegida do ambiente de execução.}
{Contexto principal: persistência local. O requisito é transversal a dados pessoais e operacionais armazenados no aplicativo.}
\reqfuncfull
{nf:respeitar_principios_de_privacidade}
{Respeitar princípios de privacidade e conformidade}
{A solução deve respeitar os princípios de privacidade e conformidade aplicáveis ao tratamento de dados pessoais no contexto brasileiro, especialmente quanto à necessidade, finalidade e restrição de uso das informações coletadas. O requisito funciona como diretriz geral de conformidade para o desenho e a evolução do sistema. Sua aplicação deve orientar tanto o tratamento dos dados quanto a exposição dessas informações aos perfis do produto.}
{Essencial}
{Depende dos requisitos [\ref{nf:tratar_dados_com_finalidade_explicita}] e [\ref{nf:restringir_visualizacao_de_dados_pessoais}]. Depende da observância das normas e princípios aplicáveis ao contexto brasileiro de tratamento de dados pessoais.}
{Contexto principal: conformidade e governança. O requisito é transversal a todo o ciclo de tratamento dos dados pessoais no produto.}

\subsubsection*{Qualidade de Uso e Evolução}
\reqfuncfull
{nf:operar_prioritariamente_offline}
{Operar prioritariamente de forma offline}
{O sistema deve operar prioritariamente de forma offline nas funcionalidades centrais de consulta, visualização e reserva previstas neste recorte. A ausência de conectividade não deve impedir o uso normal das funções principais contempladas no produto, desde que o contexto local do aplicativo esteja disponível. O requisito orienta a arquitetura e o comportamento operacional do sistema.}
{Essencial}
{Depende dos requisitos funcionais centrais de agenda, mapa, reserva, cancelamento e acompanhamento individual. Depende de armazenamento local e sincronização compatíveis com operação prioritária offline.}
{Contexto principal: arquitetura e operação. O requisito define uma condição estrutural do produto, e não apenas um comportamento opcional da interface.}
\reqfuncfull
{nf:manter_fluxo_principal_simples}
{Manter o fluxo principal de reserva simples}
{O fluxo principal de reserva deve ser simples o suficiente para que um aluno de primeira utilização consiga concluí-lo em no máximo três toques a partir do dashboard, em menos de sessenta segundos. Esse requisito estabelece critério de usabilidade e simplicidade para o principal caminho de uso do sistema. A interface deve ser desenhada para reduzir fricção e tornar a reserva rápida e intuitiva.}
{Importante}
{Depende dos requisitos funcionais do fluxo principal de agenda, seleção de turma e data, mapa e check-in. Depende de desenho de interface compatível com o limite de interações e tempo definidos.}
{Contexto principal: usabilidade do fluxo de reserva. O requisito é mensurável por número de toques e tempo de conclusão.}
\reqfuncfull
{nf:preservar_consistencia_no_contexto_mobile}
{Preservar consistência funcional e visual no contexto mobile}
{O comportamento funcional e visual das funcionalidades centrais deve permanecer consistente nos dispositivos móveis suportados pelo projeto. Diferenças técnicas entre ambientes móveis não devem alterar o entendimento do fluxo, das regras ou dos estados operacionais principais do sistema. O requisito busca preservar previsibilidade e coerência da experiência de uso no contexto principal do produto.}
{Importante}
{Depende dos requisitos funcionais centrais e do desenho de interface adotado pelo produto. Depende da implementação mobile manter equivalência de comportamento nas funcionalidades nucleares entre os dispositivos suportados.}
{Contexto principal: experiência mobile. O requisito não exige identidade absoluta de implementação, mas sim consistência percebida e funcional no ambiente principal de uso do sistema.}
\reqfuncfull
{nf:permitir_evolucao_sem_acoplamento_excessivo}
{Permitir evolução sem acoplamento excessivo}
{A solução deve ser mantida de forma extensível, permitindo evolução das entidades e fluxos do domínio sem acoplamento excessivo entre módulos. Mudanças em partes do sistema não devem impor refatorações desproporcionais em componentes não relacionados. O requisito orienta a qualidade estrutural da implementação para favorecer manutenção e crescimento controlado do produto.}
{Importante}
{Depende da organização arquitetural e modular do sistema. Depende de separação adequada entre responsabilidades do domínio, persistência, interface e regras operacionais.}
{Contexto principal: manutenibilidade e evolução técnica. O requisito é estrutural e se aplica ao desenho interno da solução.}
\reqfuncfull
{nf:preservar_confiabilidade_historica}
{Preservar confiabilidade histórica dos registros de check-in}
{O sistema deve preservar confiabilidade histórica dos registros de check-in, garantindo que cancelamentos e alterações operacionais não corrompam as métricas individuais e gerenciais derivadas do domínio. O histórico deve permanecer rastreável e coerente mesmo diante de mudanças no estado operacional das reservas. O requisito protege a utilidade analítica e a integridade temporal dos dados.}
{Importante}
{Depende dos requisitos funcionais de preservação de histórico de check-in cancelado, desconsideração de cancelados nas métricas e reflexo operacional de alterações no domínio. Depende da manutenção correta dos estados históricos dos registros.}
{Contexto principal: histórico e análise. O requisito é transversal aos fluxos que alteram estado de check-ins e métricas derivadas.}
\reqfuncfull
{nf:executar_operacoes_criticas_de_forma_atomica}
{Executar operações críticas de forma atômica}
{Operações que envolvam múltiplas escritas correlatas, como reserva, cancelamento de check-in e atualização de disponibilidade, devem ser executadas de forma atômica, sem deixar o banco em estado inconsistente. O requisito garante que alterações parciais não comprometam a integridade do domínio quando uma operação crítica for interrompida ou falhar no meio do processo. Sua aplicação é essencial para consistência operacional do sistema.}
{Essencial}
{Depende dos requisitos funcionais de reserva, cancelamento e atualização de disponibilidade. Depende de mecanismos de persistência capazes de garantir atomicidade nas operações críticas.}
{Contexto principal: persistência e consistência transacional. O requisito se aplica a toda operação com múltiplas escritas relacionadas.}
\reqfuncfull
{nf:permitir_expansao_do_padrao_cadastral}
{Permitir expansão do padrão cadastral sem refatorações invasivas}
{A inclusão de novas entidades cadastrais no padrão já adotado pelo projeto deve ser possível sem exigir refatorações invasivas nas entidades existentes. O requisito busca preservar reuso de abordagem, previsibilidade estrutural e capacidade de crescimento controlado do conjunto de cadastros do domínio. A evolução do modelo cadastral não deve degradar a estabilidade dos componentes já implementados.}
{Importante}
{Depende dos requisitos transversais de CRUD e da organização modular adotada para formulários, validações, associações e persistência das entidades. Depende de desenho técnico que favoreça extensão sem ruptura generalizada.}
{Contexto principal: evolução do ecossistema cadastral. O requisito é estrutural e orienta a expansão futura do domínio operacional.}


% --------------------------------------------------
% 6 OUTROS REQUISITOS
% --------------------------------------------------
% \newcounter{orreq}

\newcommand{\orreqid}{OR\ifnum\value{orreq}<10 00\else\ifnum\value{orreq}<100 0\fi\fi\arabic{orreq}}

\newcommand{\reqoutro}[6]{%
  \refstepcounter{orreq}%
  \subsection{[\orreqid] #2}%
  \label{#1}%
  \noindent\textbf{Descrição:} #3

  \noindent\textbf{Prioridade:} #4

  \noindent\textbf{Dependência:} #5

  \noindent\textbf{Outros:} #6

  \vspace{0.4cm}
}

\section{Outros Requisitos}

Esta seção reúne requisitos complementares que não pertencem ao núcleo funcional nem aos requisitos não funcionais já especificados nas seções anteriores. Os itens aqui descritos tratam principalmente de conformidade legal, restrições de implantação e operação, e qualidade do processo de desenvolvimento e manutenção do próprio artefato de especificação.

\subsubsection*{Conformidade Legal e Regulatória}

\reqoutro
{or:base_legal_lgpd}
{Definir base legal para tratamento de dados pessoais}
{O sistema deve apoiar a operação do produto de modo compatível com a existência de base legal explícita para o tratamento dos dados pessoais dos alunos, em conformidade com a Lei n.\ 13.709/2018. A base legal adotada para cada categoria relevante de dado pessoal deve estar definida antes da coleta e deve orientar o cadastro, o uso e a manutenção dessas informações no produto. O sistema não deve pressupor coleta irrestrita de dados pessoais sem finalidade e base legal previamente estabelecidas no contexto de uso da academia.}
{Essencial}
{Depende dos requisitos [\ref{nf:tratar_dados_com_finalidade_explicita}] e [\ref{nf:respeitar_principios_de_privacidade}]. Depende da definição, pela operação do produto, das finalidades e bases legais aplicáveis a cada categoria de dado pessoal tratada no sistema.}
{Contexto principal: conformidade com LGPD. O requisito é transversal aos fluxos que coletam, exibem, persistem ou atualizam dados pessoais dos alunos.}

\reqoutro
{or:direitos_do_titular}
{Apoiar atendimento aos direitos do titular dos dados}
{O sistema deve prever meios técnicos para que a academia identifique, localize e apresente os dados pessoais e registros relacionados a um titular quando houver solicitações ligadas aos direitos previstos na legislação aplicável. Isso inclui apoiar, quando pertinente ao contexto do produto, confirmação de tratamento, acesso, correção de dados, eliminação de dados tratados sob base adequada para isso e revogação de consentimento quando essa for a base adotada. O sistema não precisa automatizar integralmente o processo jurídico ou administrativo, mas deve fornecer condições para que a resposta ao titular seja operacionalmente viável.}
{Importante}
{Depende dos requisitos [\ref{nf:tratar_dados_com_finalidade_explicita}], [\ref{nf:restringir_visualizacao_de_dados_pessoais}] e [\ref{or:base_legal_lgpd}]. Depende da capacidade do produto de localizar os dados vinculados ao titular em seus cadastros e históricos relevantes.}
{Contexto principal: direitos do titular. O requisito prioriza apoio técnico ao atendimento das solicitações, e não a automação completa do procedimento jurídico-administrativo.}

\reqoutro
{or:retencao_de_dados}
{Definir política de retenção e descarte de dados pessoais}
{O sistema deve operar de forma compatível com uma política explícita de retenção de dados pessoais, definindo por quanto tempo cada categoria de dado deve permanecer armazenada após o encerramento do vínculo do aluno com a academia ou após evento equivalente que afete sua manutenção. Quando o prazo de retenção se encerrar e não houver fundamento para preservação adicional, os dados devem poder ser eliminados ou anonimizados de forma coerente com a política adotada. A definição dessa política deve respeitar os princípios de necessidade, finalidade e preservação histórica legítima aplicáveis ao contexto do produto.}
{Importante}
{Depende dos requisitos [\ref{or:base_legal_lgpd}], [\ref{or:direitos_do_titular}] e [\ref{nf:respeitar_principios_de_privacidade}]. Depende da definição, pela operação do produto, dos prazos e critérios de retenção aplicáveis a cada categoria de dado.}
{Contexto principal: ciclo de vida dos dados pessoais. O requisito orienta o desenho da persistência e das regras de descarte ou anonimização compatíveis com o domínio.}

\reqoutro
{or:minimizacao_de_dados}
{Restringir coleta de dados pessoais ao mínimo necessário}
{O sistema deve restringir a coleta de dados pessoais ao mínimo necessário para as finalidades declaradas do produto. Campos opcionais presentes no modelo do domínio, especialmente dados complementares de perfil ou redes sociais, não devem ser tratados como obrigatórios sem justificativa operacional clara e compatibilidade com a finalidade declarada. Sempre que houver ampliação futura do cadastro, a inclusão de novos dados pessoais deve ser precedida de revisão de finalidade, necessidade e base legal correspondente.}
{Importante}
{Depende dos requisitos [\ref{or:base_legal_lgpd}] e [\ref{nf:tratar_dados_com_finalidade_explicita}]. Depende da avaliação de necessidade de cada campo pessoal previsto no cadastro do aluno e nas demais entidades que tratem dados de pessoas.}
{Contexto principal: minimização de dados. O requisito afeta principalmente o desenho e a evolução dos formulários cadastrais e do modelo de domínio ligado a dados pessoais.}

\reqoutro
{or:restricoes_tecnicas}
{Observar restrições técnicas e acadêmicas de implantação}
{O sistema deve permanecer compatível com as restrições técnicas e acadêmicas definidas para o projeto, especialmente quanto ao uso das tecnologias escolhidas pela equipe, à possibilidade de execução em ambiente local para avaliação e à ausência de dependência obrigatória de infraestrutura paga ou indisponível no contexto acadêmico. Mudanças relevantes de stack, persistência ou arquitetura que alterem o comportamento esperado da solução devem ser refletidas nos artefatos de especificação correspondentes.}
{Importante}
{Não há dependência de requisito específico. O requisito complementa as seções de Ambiente Operacional e de Restrições de Projeto e Implementação.}
{Contexto principal: viabilidade técnica e acadêmica. O requisito orienta a manutenção de decisões tecnológicas compatíveis com a entrega e a avaliação do projeto.}

\subsubsection*{Implantação e Operação}

\reqoutro
{or:ambientes_separados}
{Permitir separação de ambientes de execução}
{A solução deve permitir configuração separada para ambientes distintos de execução, como desenvolvimento, homologação e produção, preservando independência de parâmetros, credenciais, bases de dados e demais configurações operacionais relevantes. A troca de ambiente deve ocorrer por configuração externa apropriada, sem exigir alterações manuais no código-fonte para cada publicação. Quando houver uso de dados reais, o sistema e seu processo de implantação devem evitar sua exposição em ambientes não destinados à operação efetiva.}
{Importante}
{Depende dos requisitos [\ref{or:retencao_de_dados}] e [\ref{nf:armazenar_dados_em_area_protegida}]. Depende da adoção de estratégia de configuração e persistência compatível com separação de ambientes no contexto tecnológico do projeto.}
{Contexto principal: implantação e operação. No contexto acadêmico, a granularidade exata dos ambientes pode variar, desde que exista separação suficiente para desenvolvimento e demonstração.}

\reqoutro
{or:documentacao_instalacao}
{Documentar instalação e configuração inicial}
{O produto deve ser acompanhado de documentação técnica suficiente para permitir instalação, configuração inicial e verificação básica de funcionamento sem dependência de suporte informal dos autores. Essa documentação deve cobrir, no mínimo, pré-requisitos, passos de execução ou build, configurações necessárias do ambiente, inicialização da persistência quando aplicável e forma de validar que a aplicação está operacional.}
{Importante}
{Depende do requisito [\ref{or:restricoes_tecnicas}] e da definição do processo de execução e configuração adotado pela equipe.}
{Contexto principal: entrega e operação técnica. O requisito complementa a documentação do projeto e favorece reprodutibilidade da avaliação e da manutenção futura.}

\subsubsection*{Qualidade do Processo de Desenvolvimento}

\reqoutro
{or:rastreabilidade}
{Manter rastreabilidade entre requisitos e implementação}
{A equipe deve manter rastreabilidade suficiente entre os requisitos especificados, os casos de uso, os artefatos de análise e os elementos relevantes da implementação, como telas, módulos, regras de domínio e testes quando existirem. Alterações significativas na especificação não devem reutilizar identificadores antigos de forma ambígua e devem preservar a possibilidade de compreender o impacto das mudanças sobre o produto. Essa rastreabilidade deve apoiar revisão, validação e manutenção do projeto ao longo de sua evolução.}
{Importante}
{Depende dos requisitos [\ref{or:documentacao_instalacao}] e [\ref{or:registro_de_desvios}] no que diz respeito à disciplina de manutenção documental e comunicação de mudanças. Depende também da manutenção coerente dos artefatos do projeto.}
{Contexto principal: governança do desenvolvimento. A rastreabilidade pode ser mantida por organização do repositório, referências cruzadas no documento ou artefato complementar equivalente.}

\reqoutro
{or:registro_de_desvios}
{Registrar desvios entre especificação e implementação}
{Sempre que a implementação evidenciar impossibilidade técnica, ambiguidade relevante, contradição documental ou ajuste necessário em relação a um requisito previamente especificado, a equipe deve registrar o desvio de forma explícita, indicando o requisito afetado, a natureza do problema, a decisão tomada e o impacto produzido sobre os artefatos relacionados. O objetivo é evitar divergências silenciosas entre o documento e o comportamento efetivamente implementado, preservando a integridade do processo de engenharia adotado no projeto.}
{Importante}
{Não há dependência de requisito específico. O requisito se relaciona com a manutenção do histórico de validações, com a revisão da especificação e com a governança geral do processo de desenvolvimento.}
{Contexto principal: controle de mudanças e conformidade acadêmica. O registro deve ser suficientemente objetivo para permitir auditoria posterior do motivo e do efeito da alteração realizada.}
\newcounter{orreq}

\newcommand{\orreqid}{OR\ifnum\value{orreq}<10 00\else\ifnum\value{orreq}<100 0\fi\fi\arabic{orreq}}

\newcommand{\reqoutro}[6]{%
  \refstepcounter{orreq}%
  \subsection{[\orreqid] #2}%
  \label{#1}%
  \noindent\textbf{Descrição:} #3

  \noindent\textbf{Prioridade:} #4

  \noindent\textbf{Dependência:} #5

  \noindent\textbf{Outros:} #6

  \vspace{0.4cm}
}

\section{Outros Requisitos}

Esta seção reúne requisitos complementares que não pertencem ao núcleo funcional nem aos requisitos não funcionais já especificados nas seções anteriores. Os itens aqui descritos tratam principalmente de conformidade legal, restrições de implantação e operação, e qualidade do processo de desenvolvimento e manutenção do próprio artefato de especificação.

\subsubsection*{Conformidade Legal e Regulatória}

\reqoutro
{or:base_legal_lgpd}
{Definir base legal para tratamento de dados pessoais}
{O sistema deve apoiar a operação do produto de modo compatível com a existência de base legal explícita para o tratamento dos dados pessoais dos alunos, em conformidade com a Lei n.\ 13.709/2018. A base legal adotada para cada categoria relevante de dado pessoal deve estar definida antes da coleta e deve orientar o cadastro, o uso e a manutenção dessas informações no produto. O sistema não deve pressupor coleta irrestrita de dados pessoais sem finalidade e base legal previamente estabelecidas no contexto de uso da academia.}
{Essencial}
{Depende dos requisitos [\ref{nf:tratar_dados_com_finalidade_explicita}] e [\ref{nf:respeitar_principios_de_privacidade}]. Depende da definição, pela operação do produto, das finalidades e bases legais aplicáveis a cada categoria de dado pessoal tratada no sistema.}
{Contexto principal: conformidade com LGPD. O requisito é transversal aos fluxos que coletam, exibem, persistem ou atualizam dados pessoais dos alunos.}

\reqoutro
{or:direitos_do_titular}
{Apoiar atendimento aos direitos do titular dos dados}
{O sistema deve prever meios técnicos para que a academia identifique, localize e apresente os dados pessoais e registros relacionados a um titular quando houver solicitações ligadas aos direitos previstos na legislação aplicável. Isso inclui apoiar, quando pertinente ao contexto do produto, confirmação de tratamento, acesso, correção de dados, eliminação de dados tratados sob base adequada para isso e revogação de consentimento quando essa for a base adotada. O sistema não precisa automatizar integralmente o processo jurídico ou administrativo, mas deve fornecer condições para que a resposta ao titular seja operacionalmente viável.}
{Importante}
{Depende dos requisitos [\ref{nf:tratar_dados_com_finalidade_explicita}], [\ref{nf:restringir_visualizacao_de_dados_pessoais}] e [\ref{or:base_legal_lgpd}]. Depende da capacidade do produto de localizar os dados vinculados ao titular em seus cadastros e históricos relevantes.}
{Contexto principal: direitos do titular. O requisito prioriza apoio técnico ao atendimento das solicitações, e não a automação completa do procedimento jurídico-administrativo.}

\reqoutro
{or:retencao_de_dados}
{Definir política de retenção e descarte de dados pessoais}
{O sistema deve operar de forma compatível com uma política explícita de retenção de dados pessoais, definindo por quanto tempo cada categoria de dado deve permanecer armazenada após o encerramento do vínculo do aluno com a academia ou após evento equivalente que afete sua manutenção. Quando o prazo de retenção se encerrar e não houver fundamento para preservação adicional, os dados devem poder ser eliminados ou anonimizados de forma coerente com a política adotada. A definição dessa política deve respeitar os princípios de necessidade, finalidade e preservação histórica legítima aplicáveis ao contexto do produto.}
{Importante}
{Depende dos requisitos [\ref{or:base_legal_lgpd}], [\ref{or:direitos_do_titular}] e [\ref{nf:respeitar_principios_de_privacidade}]. Depende da definição, pela operação do produto, dos prazos e critérios de retenção aplicáveis a cada categoria de dado.}
{Contexto principal: ciclo de vida dos dados pessoais. O requisito orienta o desenho da persistência e das regras de descarte ou anonimização compatíveis com o domínio.}

\reqoutro
{or:minimizacao_de_dados}
{Restringir coleta de dados pessoais ao mínimo necessário}
{O sistema deve restringir a coleta de dados pessoais ao mínimo necessário para as finalidades declaradas do produto. Campos opcionais presentes no modelo do domínio, especialmente dados complementares de perfil ou redes sociais, não devem ser tratados como obrigatórios sem justificativa operacional clara e compatibilidade com a finalidade declarada. Sempre que houver ampliação futura do cadastro, a inclusão de novos dados pessoais deve ser precedida de revisão de finalidade, necessidade e base legal correspondente.}
{Importante}
{Depende dos requisitos [\ref{or:base_legal_lgpd}] e [\ref{nf:tratar_dados_com_finalidade_explicita}]. Depende da avaliação de necessidade de cada campo pessoal previsto no cadastro do aluno e nas demais entidades que tratem dados de pessoas.}
{Contexto principal: minimização de dados. O requisito afeta principalmente o desenho e a evolução dos formulários cadastrais e do modelo de domínio ligado a dados pessoais.}

\reqoutro
{or:restricoes_tecnicas}
{Observar restrições técnicas e acadêmicas de implantação}
{O sistema deve permanecer compatível com as restrições técnicas e acadêmicas definidas para o projeto, especialmente quanto ao uso das tecnologias escolhidas pela equipe, à possibilidade de execução em ambiente local para avaliação e à ausência de dependência obrigatória de infraestrutura paga ou indisponível no contexto acadêmico. Mudanças relevantes de stack, persistência ou arquitetura que alterem o comportamento esperado da solução devem ser refletidas nos artefatos de especificação correspondentes.}
{Importante}
{Não há dependência de requisito específico. O requisito complementa as seções de Ambiente Operacional e de Restrições de Projeto e Implementação.}
{Contexto principal: viabilidade técnica e acadêmica. O requisito orienta a manutenção de decisões tecnológicas compatíveis com a entrega e a avaliação do projeto.}

\subsubsection*{Implantação e Operação}

\reqoutro
{or:ambientes_separados}
{Permitir separação de ambientes de execução}
{A solução deve permitir configuração separada para ambientes distintos de execução, como desenvolvimento, homologação e produção, preservando independência de parâmetros, credenciais, bases de dados e demais configurações operacionais relevantes. A troca de ambiente deve ocorrer por configuração externa apropriada, sem exigir alterações manuais no código-fonte para cada publicação. Quando houver uso de dados reais, o sistema e seu processo de implantação devem evitar sua exposição em ambientes não destinados à operação efetiva.}
{Importante}
{Depende dos requisitos [\ref{or:retencao_de_dados}] e [\ref{nf:armazenar_dados_em_area_protegida}]. Depende da adoção de estratégia de configuração e persistência compatível com separação de ambientes no contexto tecnológico do projeto.}
{Contexto principal: implantação e operação. No contexto acadêmico, a granularidade exata dos ambientes pode variar, desde que exista separação suficiente para desenvolvimento e demonstração.}

\reqoutro
{or:documentacao_instalacao}
{Documentar instalação e configuração inicial}
{O produto deve ser acompanhado de documentação técnica suficiente para permitir instalação, configuração inicial e verificação básica de funcionamento sem dependência de suporte informal dos autores. Essa documentação deve cobrir, no mínimo, pré-requisitos, passos de execução ou build, configurações necessárias do ambiente, inicialização da persistência quando aplicável e forma de validar que a aplicação está operacional.}
{Importante}
{Depende do requisito [\ref{or:restricoes_tecnicas}] e da definição do processo de execução e configuração adotado pela equipe.}
{Contexto principal: entrega e operação técnica. O requisito complementa a documentação do projeto e favorece reprodutibilidade da avaliação e da manutenção futura.}

\subsubsection*{Qualidade do Processo de Desenvolvimento}

\reqoutro
{or:rastreabilidade}
{Manter rastreabilidade entre requisitos e implementação}
{A equipe deve manter rastreabilidade suficiente entre os requisitos especificados, os casos de uso, os artefatos de análise e os elementos relevantes da implementação, como telas, módulos, regras de domínio e testes quando existirem. Alterações significativas na especificação não devem reutilizar identificadores antigos de forma ambígua e devem preservar a possibilidade de compreender o impacto das mudanças sobre o produto. Essa rastreabilidade deve apoiar revisão, validação e manutenção do projeto ao longo de sua evolução.}
{Importante}
{Depende dos requisitos [\ref{or:documentacao_instalacao}] e [\ref{or:registro_de_desvios}] no que diz respeito à disciplina de manutenção documental e comunicação de mudanças. Depende também da manutenção coerente dos artefatos do projeto.}
{Contexto principal: governança do desenvolvimento. A rastreabilidade pode ser mantida por organização do repositório, referências cruzadas no documento ou artefato complementar equivalente.}

\reqoutro
{or:registro_de_desvios}
{Registrar desvios entre especificação e implementação}
{Sempre que a implementação evidenciar impossibilidade técnica, ambiguidade relevante, contradição documental ou ajuste necessário em relação a um requisito previamente especificado, a equipe deve registrar o desvio de forma explícita, indicando o requisito afetado, a natureza do problema, a decisão tomada e o impacto produzido sobre os artefatos relacionados. O objetivo é evitar divergências silenciosas entre o documento e o comportamento efetivamente implementado, preservando a integridade do processo de engenharia adotado no projeto.}
{Importante}
{Não há dependência de requisito específico. O requisito se relaciona com a manutenção do histórico de validações, com a revisão da especificação e com a governança geral do processo de desenvolvimento.}
{Contexto principal: controle de mudanças e conformidade acadêmica. O registro deve ser suficientemente objetivo para permitir auditoria posterior do motivo e do efeito da alteração realizada.}
\newpage
% --------------------------------------------------
% apendice
% --------------------------------------------------
\section*{Apêndice}
\addcontentsline{toc}{section}{Apêndice}

\subsection*{Apêndice A: Levantamento de Requisitos}
\addcontentsline{toc}{subsection}{Apêndice A: Levantamento de Requisitos}
\label{ap:levantamento-requisitos}

O levantamento de requisitos considerado para a elaboração desta especificação foi estruturado de forma a refletir práticas compatíveis com projetos de software orientados à operação. O objetivo desse processo foi consolidar informações sobre o contexto da academia, os problemas centrais do fluxo de \textit{check-in}, as necessidades dos usuários e os limites do recorte funcional adotado no \ProjetoNome.

Para compor essa base de trabalho, considera-se a utilização combinada de técnicas de engenharia de requisitos, com ênfase em:

\begin{itemize}[leftmargin=1.2cm]
    \item \textbf{Entrevistas semiestruturadas}: utilizadas para levantar dores operacionais, expectativas de uso e percepções sobre os principais problemas ligados à reserva, à ocupação das aulas, à manutenção das bikes e à experiência do aluno.
    \item \textbf{Sessões colaborativas do tipo \textit{JAD}}: utilizadas para consolidar entendimento entre as partes interessadas sobre o fluxo central do produto, suas prioridades e os limites do escopo.
    \item \textbf{Análise documental}: utilizada para revisar artefatos do próprio projeto, regras de negócio, diagrama de classes, protótipos de interface e materiais de apoio relevantes para a especificação.
    \item \textbf{Consolidação de artefatos}: utilizada para transformar as informações coletadas em um conjunto coerente de definições, escopo, regras e requisitos.
\end{itemize}

No contexto desta especificação, essas técnicas sustentam a identificação do \textit{check-in} como núcleo do produto, da Professora como frente de suporte operacional e do grupo de alunos como estrutura complementar de menor prioridade. Também sustentam a delimitação entre o que faz parte do recorte atual do sistema e o que deve permanecer como possibilidade de evolução futura. Em conjunto, esse processo permitiu consolidar uma especificação alinhada ao problema de negócio, ao domínio modelado e às decisões de escopo formalizadas neste documento.

\newpage

\subsection*{Apêndice B: Modelos de Análise}
\addcontentsline{toc}{subsection}{Apêndice B: Modelos de Análise}

\subsubsection*{B.1 Diagrama de Classes}
\begin{figure}[!h]
\centering
\includegraphics[width=\textwidth]{img/dc.png}
\caption{Diagrama de classes do projeto \ProjetoNome. Para preservar clareza visual, foram omitidos os fluxos e operacoes tradicionais de CRUD das entidades de cadastro.}
\end{figure}
\newpage

\subsubsection*{B.2 Diagramas de Casos de Uso}
\paragraph{Perfil Professora}
\begin{figure}[!h]
\centering
\includegraphics[width=0.9\textwidth]{img/dcup.png}
\caption{Diagrama de casos de uso do perfil Professora no projeto \ProjetoNome.}
\end{figure}

\newpage

\paragraph{Perfil Aluno}
\begin{figure}[!h]
\centering
\includegraphics[width=0.9\textwidth]{img/dcug_aluno.png}
\caption{Diagrama de casos de uso do perfil Aluno no projeto \ProjetoNome.}
\end{figure}
\newpage

\subsubsection*{B.3 Tabela de Casos de Uso}
\begin{longtable}{@{}>{\raggedright\arraybackslash}p{0.10\textwidth}>{\raggedright\arraybackslash}p{0.52\textwidth}>{\centering\arraybackslash}p{0.14\textwidth}>{\centering\arraybackslash}p{0.14\textwidth}@{}}
\toprule
\textbf{ID} & \textbf{Caso de Uso} & \textbf{Professora} & \textbf{Aluno} \\
\midrule
\endfirsthead
\toprule
\textbf{ID} & \textbf{Caso de Uso} & \textbf{Professora} & \textbf{Aluno} \\
\midrule
\endhead
UC01 & Autenticar no Sistema & Sim & Sim \\
UC02 & Realizar Checkin & Sim & Sim \\
UC03 & Verificar Disponibilidade da Bike (\textit{include} de UC02) & --- & --- \\
UC04 & Visualizar Mapa da Sala (\textit{include} de UC02) & --- & --- \\
UC05 & Cancelar Checkin (\textit{extend} de UC02) & Sim, qualquer & Sim, proprio \\
UC06 & Gerenciar Alunos e Grupos & Sim & --- \\
UC07 & Gerenciar Bikes e Fabricantes & Sim & --- \\
UC08 & Gerenciar Manutencoes & Sim & --- \\
UC09 & Cancelar Manutencao (\textit{extend} de UC08) & Sim & --- \\
UC10 & Gerenciar Salas e Turmas & Sim & --- \\
UC11 & Verificar Disponibilidade de Horario (\textit{include} de UC10) & --- & --- \\
UC12 & Associar Mix a Turma (\textit{extend} de UC10) & Sim & --- \\
UC13 & Gerenciar Mixes e Repertorio & Sim & --- \\
UC14 & Consultar Agenda Semanal & Sim & Sim \\
UC15 & Consultar Mix e Repertorio & Sim & Sim \\
UC16 & Consultar Historico de Presenca & --- & Sim \\
UC17 & Visualizar Dashboard & Sim & Sim \\
UC18 & Gerar Relatorios Gerenciais & Sim & --- \\
\bottomrule
\end{longtable}

\subsubsection*{B.4 Descricoes dos Casos de Uso}
\begin{description}[style=nextline,leftmargin=0pt,labelsep=0pt]
\item[UC01 -- Autenticar no Sistema] O usuario informa suas credenciais para acessar o sistema. O perfil autenticado determina quais funcionalidades estarao disponiveis apos o login. A senha deve possuir no minimo 6 caracteres.
\item[UC02 -- Realizar Checkin] O aluno ou a professora reserva uma posicao de bike em uma turma para uma data especifica. O sistema valida a disponibilidade da bike e exibe o mapa da sala antes de confirmar a reserva. A data deve corresponder a um dos dias da semana da turma e nao pode ser anterior a data atual.
\item[UC03 -- Verificar Disponibilidade da Bike (\textit{include})] Invocado sempre que UC02 e executado. Verifica se a posicao esta dentro dos limites da sala, se nao ha checkin ativo na mesma posicao, turma e data, e se a turma ainda possui vaga disponivel.
\item[UC04 -- Visualizar Mapa da Sala (\textit{include})] Invocado sempre que UC02 e executado. Exibe a grade visual da sala com o estado de cada posicao: livre, reservada, em manutencao ou reservada para a professora. A professora visualiza o nome do aluno em posicoes ocupadas; o aluno ve apenas ocupacao anonima.
\item[UC05 -- Cancelar Checkin (\textit{extend})] Estende UC02 quando existe um checkin passivel de cancelamento. Libera a posicao, atualiza a disponibilidade da turma e mantem o historico com cancelamento registrado. O aluno cancela apenas o proprio checkin enquanto a data nao passou; a professora pode cancelar qualquer checkin.
\item[UC06 -- Gerenciar Alunos e Grupos] Realiza o CRUD de `Aluno` e `GrupoAlunos`, incluindo ativacao e inativacao logica. A inativacao de aluno cancela seus checkins futuros. Grupos com todos os alunos inativos devem ser inativados automaticamente.
\item[UC07 -- Gerenciar Bikes e Fabricantes] Realiza o CRUD de `Bike` e `Fabricante`. Toda bike deve estar associada a um fabricante. A inativacao de bike cancela checkins futuros; a inativacao do fabricante nao inativa automaticamente as bikes associadas.
\item[UC08 -- Gerenciar Manutencoes] Registra ocorrencias de manutencao para uma bike especifica, com tipo e datas. Ao abrir uma manutencao, a bike fica indisponivel. Ao registrar a realizacao, a bike volta a ficar disponivel. Uma bike so pode ter uma manutencao pendente por vez.
\item[UC09 -- Cancelar Manutencao (\textit{extend})] Estende UC08 quando existe manutencao passivel de cancelamento. Define o cancelamento, recalcula a disponibilidade da bike e preserva o historico da ocorrencia.
\item[UC10 -- Gerenciar Salas e Turmas] Realiza o CRUD de `Sala` e `Turma`. A sala define o layout fisico e as posicoes de bikes. A turma associa um horario fixo a uma sala. Ao ativar uma turma, o sistema deve verificar a disponibilidade do horario na sala.
\item[UC11 -- Verificar Disponibilidade de Horario (\textit{include})] Invocado quando UC10 ativa ou edita uma turma. Verifica se nao ha sobreposicao de dias da semana e horario com outra turma ativa na mesma sala.
\item[UC12 -- Associar Mix a Turma (\textit{extend})] Estende UC10 de forma opcional. Cria o vinculo `TurmaMix` ativo para a turma. Quando um novo mix e associado, o vinculo anterior deve ser encerrado com `dataFim`. Uma turma pode permanecer ativa sem mix associado.
\item[UC13 -- Gerenciar Mixes e Repertorio] Realiza o CRUD de `Mix`, `Musica`, `ArtistaBanda`, `CategoriaMusica` e `VideoAula`. Um mix deve ter ao menos uma musica. A mesma musica nao pode se repetir no mesmo mix. Inativar um mix encerra o vinculo ativo com a turma.
\item[UC14 -- Consultar Agenda Semanal] Exibe as turmas ativas em grade semanal, com sala, mix atual e quantidade de bikes disponiveis em tempo real. Deve permitir filtros por sala ou dia e indicar visualmente turmas com alta ocupacao.
\item[UC15 -- Consultar Mix e Repertorio] Exibe o mix ativo da turma, a lista de musicas com artista e categorias, links de video-aula e o historico recente de mixes utilizados.
\item[UC16 -- Consultar Historico de Presenca] Exibe ao aluno seu historico de checkins, com total de aulas concluidas, recortes anual e mensal, sequencia de presenca e posicoes favoritas por sala.
\item[UC17 -- Visualizar Dashboard] Para a professora, apresenta KPIs operacionais, como alunos ativos, aulas agendadas, mixes em uso e bikes disponiveis. Para o aluno, apresenta metricas individuais de aulas concluidas e aulas agendadas.
\item[UC18 -- Gerar Relatorios Gerenciais] Disponibiliza relatorios estrategicos e gerenciais relacionados a retencao, saude da frota, ocupacao das turmas, linha do tempo musical e frequencia por turma.
\end{description}

\subsubsection*{B.5 Legenda dos Relacionamentos}
\begin{tabularx}{\textwidth}{@{}>{\raggedright\arraybackslash}p{0.28\textwidth}>{\raggedright\arraybackslash}p{0.16\textwidth}>{\raggedright\arraybackslash}X@{}}
\toprule
\textbf{Notacao} & \textbf{Tipo} & \textbf{Semantica} \\
\midrule
\texttt{<<include>>} & Include & O caso de uso base sempre invoca o caso de uso incluido; trata-se de fluxo obrigatorio. \\
\texttt{<<extend>>} & Extend & O caso de uso de extensao adiciona comportamento opcional ao caso de uso base, de forma condicional. \\
Associacao & Associacao & O ator participa diretamente do caso de uso. \\
\bottomrule
\end{tabularx}

\newpage

\subsection*{Apêndice C: Prototipações das Interfaces}
\addcontentsline{toc}{subsection}{Apêndice C: Prototipações das Interfaces}

\subsubsection*{C.1 Dashboard da Professora}
\begin{figure}[!h]
\centering
\includegraphics[width=0.42\textwidth]{img/widgets/01-dashboard_prof.png}
\caption{Prototipo do dashboard da professora.}
\end{figure}
\noindent Esta prototipacao apresenta a visao inicial do perfil da professora, com indicadores operacionais e navegacao para as areas principais do sistema. A tela funciona como ponto de entrada para os fluxos de suporte operacional ligados a turmas, bikes, salas e manutencoes.

\noindent\textbf{Requisitos relacionados:} [\ref{rf:manter_salas}], [\ref{rf:manter_bikes}], [\ref{rf:manter_turmas}], [\ref{rf:registrar_manutencoes}] e [\ref{rf:refletir_efeitos_de_manutencao}].

\newpage

\subsubsection*{C.2 Dashboard da Professora -- Area de Cadastros}
\begin{figure}[!h]
\centering
\includegraphics[width=0.42\textwidth]{img/widgets/02-dashboard_prof_cadastros.png}
\caption{Prototipo da area de cadastros acessivel pelo dashboard da professora.}
\end{figure}
\noindent Esta tela evidencia a organizacao dos acessos de cadastro e manutencao das entidades operacionais do dominio. Seu papel e centralizar os fluxos de criacao e edicao, reduzindo friccao para a professora na administracao cotidiana do sistema.

\noindent\textbf{Requisitos relacionados:} [\ref{rf:validar_formularios_cadastrais}], [\ref{rf:alteracao_controlada_cadastral}], [\ref{rf:manter_salas}], [\ref{rf:manter_bikes}] e [\ref{rf:manter_turmas}].

\newpage

\subsubsection*{C.3 Dashboard da Professora -- Area de Listas}
\begin{figure}[!h]
\centering
\includegraphics[width=0.42\textwidth]{img/widgets/03-dashboard_prof_listas.png}
\caption{Prototipo da area de listagens acessivel pelo dashboard da professora.}
\end{figure}
\noindent Esta prototipacao mostra a navegacao para consultas operacionais e listas administrativas. Ela apoia os fluxos de selecao de registros para edicao, inativacao logica e acompanhamento das entidades mantidas pela professora.

\noindent\textbf{Requisitos relacionados:} [\ref{rf:exclusao_logica_cadastral}], [\ref{rf:alteracao_controlada_cadastral}], [\ref{rf:manter_salas}], [\ref{rf:manter_bikes}] e [\ref{rf:manter_turmas}].

\newpage

\subsubsection*{C.4 Dashboard da Professora -- Area de Aulas}
\begin{figure}[!h]
\centering
\includegraphics[width=0.42\textwidth]{img/widgets/04-dashboard_prof_aulas.png}
\caption{Prototipo da area de aulas acessivel pelo dashboard da professora.}
\end{figure}
\noindent Esta tela concentra acessos relacionados ao planejamento e acompanhamento das aulas. A organizacao proposta dialoga com os requisitos de manutencao de turmas e com a visualizacao operacional da ocupacao associada a cada aula.

\noindent\textbf{Requisitos relacionados:} [\ref{rf:manter_turmas}], [\ref{rf:validar_ativacao_de_turma}], [\ref{rf:visualizar_mapa_operacional}] e [\ref{rf:cancelar_checkins_de_alunos}].

\newpage

\subsubsection*{C.5 Dashboard da Professora -- Area de Manutencao}
\begin{figure}[!h]
\centering
\includegraphics[width=0.42\textwidth]{img/widgets/05-dashboard_prof_manutencao.png}
\caption{Prototipo da area de manutencao acessivel pelo dashboard da professora.}
\end{figure}
\noindent Esta prototipacao evidencia o agrupamento dos fluxos de manutencao em uma area especifica da interface. A proposta e coerente com o uso frequente da professora para registrar ocorrencias e acompanhar seus efeitos operacionais sobre as bikes e as turmas.

\noindent\textbf{Requisitos relacionados:} [\ref{rf:registrar_manutencoes}] e [\ref{rf:refletir_efeitos_de_manutencao}].

\newpage

\subsubsection*{C.6 Cadastro de Fabricante}
\begin{figure}[!h]
\centering
\includegraphics[width=0.42\textwidth]{img/widgets/06-cadastro_fabricante.png}
\caption{Prototipo do formulario de cadastro de fabricante.}
\end{figure}
\noindent Esta tela demonstra um formulario cadastral com agrupamento de informacoes, validacao de campo obrigatorio e controle de ativacao logica. Embora o corpo principal da especificacao nao detalhe um requisito exclusivo para fabricante, a prototipacao atende ao padrao transversal de formularios administrativos do sistema.

\noindent\textbf{Requisitos relacionados:} [\ref{rf:validar_formularios_cadastrais}], [\ref{rf:alteracao_controlada_cadastral}] e [\ref{rf:exclusao_logica_cadastral}].

\newpage

\subsubsection*{C.7 Cadastro de Sala}
\begin{figure}[!h]
\centering
\includegraphics[width=0.42\textwidth]{img/widgets/07-cadastro_sala.png}
\caption{Prototipo do formulario de cadastro de sala.}
\end{figure}
\noindent Esta prototipacao materializa o cadastro da entidade sala, incluindo informacoes estruturais necessarias para a definicao da grade fisica do ambiente. A tela apoia diretamente a configuracao do layout que sustenta o mapa operacional das aulas e as reservas por posicao.

\noindent\textbf{Requisitos relacionados:} [\ref{rf:manter_salas}], [\ref{rf:definir_grade_da_sala}], [\ref{rf:validar_formularios_cadastrais}] e [\ref{rf:alteracao_controlada_cadastral}].

\newpage

\subsubsection*{C.8 Cadastro de Bike}
\begin{figure}[!h]
\centering
\includegraphics[width=0.42\textwidth]{img/widgets/08-cadastro_bike.png}
\caption{Prototipo do formulario de cadastro de bike.}
\end{figure}
\noindent Esta tela apresenta o cadastro da bike como entidade operacional associada a outras informacoes do dominio. O formulario reforca a necessidade de validacao dos campos e de uso de associacoes controladas para preservar consistencia entre bike, fabricante e uso operacional na sala.

\noindent\textbf{Requisitos relacionados:} [\ref{rf:manter_bikes}], [\ref{rf:associacao_controlada_entre_entidades}], [\ref{rf:validar_formularios_cadastrais}] e [\ref{rf:alteracao_controlada_cadastral}].

\newpage

\subsubsection*{C.9 Exemplo de Campo de Data}
\begin{figure}[!h]
\centering
\includegraphics[width=0.42\textwidth]{img/widgets/09-cadastro_bike_campo_data.png}
\caption{Prototipo destacando o uso de campo de data em formulario cadastral.}
\end{figure}
\noindent Esta imagem evidencia o tratamento de campos de data na interface, com foco em padronizacao de entrada e prevencao de erros de preenchimento. O componente demonstra como a interface contribui para o cumprimento das validacoes semanticas exigidas pelos formularios do sistema.

\noindent\textbf{Requisitos relacionados:} [\ref{rf:validar_formularios_cadastrais}] e [\ref{rf:manter_bikes}].

\newpage

\subsubsection*{C.10 Exemplo de Campo de Selecao Unica}
\begin{figure}[!h]
\centering
\includegraphics[width=0.42\textwidth]{img/widgets/09-cadastro_bike_campo_selecao_unico.png}
\caption{Prototipo destacando o uso de campo de selecao unica em formulario cadastral.}
\end{figure}
\noindent Esta prototipacao ilustra um campo relacional de escolha unica, adequado para associacoes entre entidades previamente cadastradas. O componente reforca a regra de que relacionamentos do dominio devem ser tratados por selecao controlada e nao por digitacao livre.

\noindent\textbf{Requisitos relacionados:} [\ref{rf:associacao_controlada_entre_entidades}], [\ref{rf:validar_formularios_cadastrais}] e [\ref{rf:manter_bikes}].

\newpage

\subsubsection*{C.11 Cadastro de Mix}
\begin{figure}[!h]
\centering
\includegraphics[width=0.42\textwidth]{img/widgets/10-cadastro_mix.png}
\caption{Prototipo do formulario de cadastro de mix.}
\end{figure}
\noindent Esta tela demonstra o formulario de definicao de um mix, com nome, periodo de uso, estado ativo e composicao do repertorio. A prototipacao e particularmente relevante para a coerencia entre vigencia, repertorio associado e futura exibicao do mix ao aluno.

\noindent\textbf{Requisitos relacionados:} [\ref{rf:associacao_controlada_entre_entidades}], [\ref{rf:validar_formularios_cadastrais}], [\ref{rf:determinar_mix_ativo}], [\ref{rf:exibir_nome_e_periodo_do_mix}] e [\ref{rf:preservar_consistencia_da_vigencia_do_mix}].

\newpage

\subsubsection*{C.12 Campo de Selecao Multipla no Cadastro de Mix}
\begin{figure}[!h]
\centering
\includegraphics[width=0.52\textwidth]{img/widgets/10-cadastro_mix_selecao_múltipla.png}
\caption{Prototipo anotado do componente de selecao multipla utilizado no cadastro de mix.}
\end{figure}
\noindent Esta imagem destaca o comportamento esperado do componente de selecao multipla para associacao de musicas ao mix. A anotacao evidencia as acoes de adicionar e remover itens selecionados, bem como a visualizacao das escolhas ja realizadas, caracteristicas diretamente alinhadas ao padrao de associacao controlada entre entidades.

\noindent\textbf{Requisitos relacionados:} [\ref{rf:associacao_controlada_entre_entidades}], [\ref{rf:exibir_musicas_do_mix}], [\ref{rf:exibir_dados_musicais}] e [\ref{rf:apresentar_links_de_video_aula}].



\end{document}